%=====================================================================
%  CONCENTRICITY OVER THE OCTONIONS -- MASTER
%  Part 1: classical zeta facts (continuation, functional equation,
%          infinite Euler product, infinite Hadamard product) -- cited.
%  Part 2: slice-preserving theory, the ring R under the *-product, and the
%          equivalence/translation theorems (concentric zero-spheres <=> RH).
%  Part 3: the categorical construction and the concentricity question.
%          The base is B = PGL(2,R) \ltimes OnePoint(R) and the sphere world
%          SphereWorld. Conditions C1--C3 determine one distinguished
%          Mobius/exponential element; orbit--stabilizer carries its whole
%          O*-valued action across the base; the value-bearing functor
%          A_A : B -> Grpd has normalized A-value states as fibres, and
%          T_A = Int_B A_A is its total. The C-residue inverse-image total
%          is transitive, hence connected, with singleton pi_0 and singleton
%          component colimit (proved, kernel-clean). That singleton does NOT
%          identify the real reads of the residue representatives: the
%          unrestricted concentricity statement is FALSE (a field-complete
%          nonconcentric A-section is exhibited), and for zeta_O the read
%          equality is equivalent to RH. SYNCHRONIZED 2026-09-02.
%=====================================================================

\documentclass[11pt]{article}

\usepackage[utf8]{inputenc}
\usepackage[margin=1in]{geometry}
\usepackage{amsmath,amssymb,amsthm,mathtools}
\usepackage{stmaryrd}   % \fatsemi — diagrammatic composition
\usepackage{amscd}
\usepackage{tikz-cd}
\usepackage{enumitem}
\usepackage{xcolor}
\usepackage[colorlinks=true,linkcolor=blue,citecolor=blue,urlcolor=blue]{hyperref}

% Vertical spacing.  Without \raggedbottom, a page that cannot fit an
% unbreakable display or tikzcd stretches the glue between its paragraphs to
% fill the height, which reads as a gap in the middle of the text.  This sends
% the slack to the foot of the page instead.
\raggedbottom
% Let displays sit closer to a page break before being pushed to the next page.
\predisplaypenalty=0
\postdisplaypenalty=0

%---------------------------------------------------------------------
%  Macros
%---------------------------------------------------------------------
\newcommand{\OO}{\mathbb{O}}
\newcommand{\CC}{\mathbb{C}}
\newcommand{\RR}{\mathbb{R}}
\newcommand{\HH}{\mathbb{H}}
\newcommand{\ZZ}{\mathbb{Z}}
\newcommand{\QQ}{\mathbb{Q}}
\newcommand{\NN}{\mathbb{N}}
\newcommand{\Ostar}{\OO^{*}}
\newcommand{\Cstar}{\CC^{*}}
\newcommand{\zetaO}{\zeta_{\Ostar}}
\newcommand{\zetaC}{\zeta_{\Cstar}}
\DeclareMathOperator{\Spec}{Spec}
\DeclareMathOperator{\Hom}{Hom}
\DeclareMathOperator{\Aut}{Aut}
\DeclareMathOperator{\re}{Re}
\DeclareMathOperator{\im}{Im}
\providecommand{\Gtwo}{G_2}

% Conspicuous markers (visible in compiled output):
\newcommand{\TODO}[1]{\medskip\noindent\textcolor{red}{\textbf{[TODO:} #1\textbf{]}}\medskip}
\newcommand{\PROGRAM}[1]{\medskip\noindent\textcolor{red}{\textbf{[PROGRAM:} #1\textbf{]}}\medskip}
\newcommand{\SOURCE}[1]{\textcolor{purple}{\small$\langle$source: #1$\rangle$}}
\newcommand{\stag}[1]{\href{https://stacks.math.columbia.edu/tag/#1}{\texttt{#1}}}
\providecommand{\Grpd}{\mathrm{Grpd}}
\providecommand{\SphereWorld}{\mathrm{SphereWorld}}
\providecommand{\nfr}{\mathfrak{n}}

% Lean blueprint macros: no-ops when this file is compiled standalone; the
% leanblueprint package redefines them when the dependency graph and the Lean
% linkage are built. Keeping them here lets every statement carry its Lean
% declaration and its dependencies without any Lean names appearing in the prose.
\providecommand{\lean}[1]{}
\providecommand{\leanok}{}
\providecommand{\authorobject}[2]{}
\providecommand{\notready}{}
\providecommand{\mathlibok}{}
\providecommand{\uses}[1]{}
\providecommand{\proves}[1]{}

%---------------------------------------------------------------------
%  Theorem environments (document-wide numbering)
%---------------------------------------------------------------------
\theoremstyle{plain}
\newtheorem{theorem}{Theorem}[section]
\newtheorem{proposition}[theorem]{Proposition}
\newtheorem{lemma}[theorem]{Lemma}
\newtheorem{corollary}[theorem]{Corollary}
\theoremstyle{definition}
\newtheorem{definition}[theorem]{Definition}
\newtheorem{condition}[theorem]{Condition}
\theoremstyle{remark}
\newtheorem{remark}[theorem]{Remark}
\newtheorem{example}[theorem]{Example}

\title{Concentric Zero-Spheres of a class of sections of the commutative ring $\mathcal R$\\ of slice preserving functions
over the compactified Octonions\\[0.4em]
\large
with the Riemann Hypothesis as a concentricity statement}
\author{Jesse Michael Paul\\[2pt]
{\normalsize PhD Candidate, Department of Mathematics \& Statistics}\\
{\normalsize University of North Carolina at Greensboro}}
\hypersetup{pdfauthor={Jesse Michael Paul}}
\date{June 2026}

\begin{document}
\maketitle

\begin{quote}\small\itshape
``\dots a `viewpoint' by itself remains fragmentary. It reveals to us one of the aspects of a
scenery or panorama, among a multiplicity of others which are equally valuable, equally `real'.
It is when complementary viewpoints of a common reality are conjugated, that is, when our `eyes'
are multiplied, that the gaze is able to penetrate further ahead in the reckoning of things.
\dots it also happens, sometimes, that a sheaf of viewpoints converging to a unique and vast
scenery gives rise to a novel thing; a thing which transcends each of the partial perspectives,
in the same way that a living being transcends each of its limbs and organs. This new thing could
be called a vision.''
\begin{flushright}\upshape --- A.\ Grothendieck, \emph{R\'ecoltes et Semailles}\end{flushright}
\end{quote}



\begin{abstract}
Let $\mathcal R$ be the commutative ring of slice-preserving slice-regular functions on the
compactified octonions $\OO^{*}=S^8$, under the regular ($\ast$) product. Part~1 gathers the
classical facts about the Riemann zeta function. Part~2 develops the slice-preserving theory,
builds $\mathcal R$, places the octonionic zeta $\zetaO$ inside it, and proves the zero and
equivalence theorems: the classically nontrivial zeros of $\zeta$ are residue-$\CC$ zero
$6$-spheres, and the Riemann Hypothesis holds exactly when those spheres are concentric, sharing
a single real centre.

Part~3 studies the \emph{concentricity question} for a section $A\in\mathcal R$
satisfying four conditions: \textup{(C1)} meromorphic continuation with a single simple pole
carried to $\infty=N$; \textup{(C2)} an infinite Euler product; \textup{(C3)} an infinite
slice-regular Weierstrass factorization whose spherical factors are the residue-$\CC$ zeros; and
\textup{(C4)} infinitely many such zeros. The question is whether the residue-$\CC$ zero
$6$-spheres of such an $A$ are concentric.

The tool is a concentricity detector built from the section itself, out of groupoids,
orbit--stabilizer, and categorical homotopy theory. Conditions \textup{(C1)--(C3)} determine one
distinguished exponential M\"obius action on the projective base
$\mathcal B=PGL(2,\RR)\ltimes\operatorname{OnePoint}(\RR)$: a logarithmic coordinate $z$ is
exponentiated in the diagonal matrix $\operatorname{diag}(e^z,1)$ and conjugated into the Cayley
disk chart, with the Euler product presenting the action at $0$ and the Weierstrass product at
$N$. The slice-preserving normalization is $G_2$-equivariant on $\Ostar$, and
orbit--stabilizer carries this whole point-valued action across the base. Its sphere-world
projection records the directions and M\"obius legs, while the genuine A-section functor
$\mathcal A_A:\mathcal B\to\mathbf{Grpd}$ has the resulting normalized A-value states and their
transports as its fibres. Its Grothendieck construction is
$\mathcal T_A=\int_{\mathcal B}\mathcal A_A$.

The C-residue system is the inverse-image action groupoid
$\mathcal R_A$ selected inside $\mathcal A_A$, with fully faithful inclusion
$\iota_A:\mathcal R_A\Rightarrow\mathcal A_A$, and hence a fully faithful inclusion of totals
$\int\iota_A:\int_{\mathcal B}\mathcal R_A\hookrightarrow\mathcal T_A$. The relative
orbit--stabilizer comparison makes this inverse-image total transitive, therefore connected, and
its component colimit is consequently a singleton. Each certified residue representative carries
a real read, the real coordinate of its zero-sphere. The singleton identifies the component
classes of any two representatives; it does not identify their reads, because no cocone,
descent law, or invariance law for the read along the transports is established, and none can
be in general: the class of A-sections contains $(q^{2}+1)\,\zetaO$, a member satisfying
\textup{(C1)--(C4)} whose zero-spheres are centred at $0$ and at the real parts of the zeta
zeros. The unrestricted concentricity theorem asserted in earlier versions of this paper is
therefore false, and is withdrawn.

What survives is exact and zeta-specific. Pairwise equality of the certified reads is the same
proposition as concentricity of the enumerated zero-spheres, and for $\zetaO$ both are
equivalent to the Riemann Hypothesis; under the Riemann Hypothesis the common centre and the
common read are $\tfrac12$. These are geometric and categorical reformulations of the Riemann
Hypothesis, formalized in Lean without placeholders. They are not a proof of it.
\end{abstract}


\tableofcontents

%=====================================================================
\section*{Introduction}
%=====================================================================
The Riemann Hypothesis is usually met head-on, as a question about where the zeros of a single complex function fall along a line. This paper approaches it from the side. It asks instead whether a family of geometric objects, sitting a few dimensions up in the octonions, share a common centre. For the octonionic zeta that question is the Riemann Hypothesis itself, seen from the side: not a route around it, but the same fact cast a few dimensions up, where more of its structure is visible.

A nontrivial zero of the zeta function, seen from the complex line alone, is a lone point in a strip. Seen from the compactified octonions, that same zero opens into a six-dimensional sphere, and the critical-line question becomes a question of arrangement: are these spheres concentric, or scattered into different hyperplanes? It takes the whole continuum of imaginary directions in the octonions, read together, to see whether the alignment holds.

The paper is built to make that convergence visible and, at every step, checkable. Part 1 recalls the classical facts about the Riemann zeta function we will lean on. Part 2 develops the slice-preserving theory of functions on the octonions, builds the commutative ring $\mathcal R$ in which the octonionic zeta lives, and proves the reframing at the centre of everything: the classically nontrivial zeros are exactly the residue-$\CC$ zero six-spheres, and the Riemann Hypothesis holds precisely when those spheres are concentric. Part 3 builds a small piece of machinery I think of as a concentricity detector, and reports exactly what it detects. From the section itself we build one distinguished action on a projective base, carry it across a continuum of Riemann spheres by orbit and stabilizer, and read off, through categorical homotopy theory, that the zero-sphere states land in a single connected piece. That is proved. What it does not give is agreement of the real reads carried by those states: the class of sections considered contains a member whose zero-spheres are not concentric, so no argument at this generality can. For the octonionic zeta, agreement of the reads is equivalent to the Riemann Hypothesis, and under the Riemann Hypothesis the common centre is one half.

No one can picture an eight-dimensional graph in their head, but tools from categorical homotopy theory, like the total-object Grothendieck construction and standard colimit arguments, let us probe this structure and reason about it. The argument is also formalized in Lean, so that each step can be checked; the reader can verify each step, and see exactly where the connected component ends and the Riemann Hypothesis begins.

%=====================================================================
\part{Classical Background}
%=====================================================================
\section{The Riemann zeta function}

The Riemann zeta function is classical, and we take its basic theory as given. We recall only what
the construction uses: the meromorphic continuation and functional equation, the Euler product over
the primes, the Hadamard product over the zeros, and Hardy's theorem. We then read $\zeta$ in
compactified form, as a map to the Riemann sphere, which is the form in which it will extend to the
octonions.

\begin{theorem}[Meromorphic continuation and functional equation; Riemann \cite{Riemann1859}]\label{thm:riemann}
\lean{riemannZeta_meromorphicOn, riemannZeta_orderAt_one,
riemannZeta_conj, riemannZeta_one_sub_zero, nontrivialZero_re_mem_Ioo}
The Dirichlet series $\zeta(s)=\sum_{n\ge1}n^{-s}$, convergent for $\re s>1$, continues to a
meromorphic function on $\CC$ with a single simple pole, at $s=1$; equivalently, a holomorphic map
to the Riemann sphere
\[
  \zeta:\Cstar\longrightarrow\Cstar,\qquad \Cstar=\CC\cup\{\infty\}=S^2,\qquad \zeta(1)=\infty .
\]
The completed function $\xi(s)=\tfrac12 s(s-1)\pi^{-s/2}\Gamma(s/2)\zeta(s)$ is entire and satisfies
\[
  \xi(s)=\xi(1-s),\qquad \xi(\bar s)=\overline{\xi(s)} .
\]
Consequently, if $\rho$ is a nontrivial zero of $\zeta$ then so are $\bar\rho$, $1-\rho$, and
$1-\bar\rho$. The trivial zeros of $\zeta$ are at $s=-2,-4,\dots$; the nontrivial zeros lie in the
strip $0<\re s<1$. The Riemann Hypothesis asserts that every nontrivial zero has $\re s=\tfrac12$.
\end{theorem}

\begin{theorem}[Infinite Euler product; {\cite[Ch.~1]{Titchmarsh86}}]\label{thm:euler}
\lean{zetaEulerLog, zetaC2_summable, zetaC2_zero_free, zetaC2_euler}
For $\re s>1$,
\[
  \zeta(s)=\prod_{p\ \mathrm{prime}}\bigl(1-p^{-s}\bigr)^{-1}
         =\exp\!\Bigl(\sum_{p}\sum_{k\ge1}\tfrac1k\,p^{-ks}\Bigr),
\]
the product taken over the infinitely many primes; it converges absolutely and is zero-free on
$\re s>1$.
\end{theorem}

\begin{theorem}[Infinite Hadamard product; {\cite[Ch.~2]{Titchmarsh86}}]\label{thm:hadamard}
\lean{xi_entire, xi_zeros_eq_nontrivialZeros, zeta_hadamard,
riemannZeta_nontrivialZeros_infinite}
$\xi$ is entire of order $1$ with infinitely many zeros, which are exactly the nontrivial zeros
$\{\rho\}$ of $\zeta$, and admits the Hadamard factorization
\[
  \xi(s)=\xi(0)\prod_{\rho}\Bigl(1-\tfrac{s}{\rho}\Bigr)e^{s/\rho},
\]
the product taken over the infinitely many nontrivial zeros.
\end{theorem}

\begin{corollary}[Infinitude of the nontrivial zeros]\label{cor:hadamard-infinitude}
\lean{riemannZeta_nontrivialZeros_infinite}
\uses{thm:hadamard}
$\zeta$ has infinitely many nontrivial zeros.
\end{corollary}
\begin{proof}
Immediate from Theorem~\ref{thm:hadamard}: its zeros are infinitely many and are exactly
the nontrivial zeros of $\zeta$.
\end{proof}

\begin{theorem}[Hardy \cite{Hardy14}]\label{thm:hardy}
Infinitely many nontrivial zeros of $\zeta$ lie on the line $\re s=\tfrac12$.
\end{theorem}

\subsection*{The compactified Riemann zeta on \texorpdfstring{$\Cstar$}{C*}}

For the slice construction, $\zeta$ is taken in compactified form, with values on the
Riemann sphere $\Cstar=\CC\cup\{\infty\}=S^2$ (the one-point compactification of $\CC$
\cite{Munkres00}). The classical facts above are retained as stated.

\begin{definition}[Compactified classical zeta]\label{def:zeta-Cstar}
\lean{zetaC, zetaC_coe, zetaC_one, zetaC_infty}
\uses{thm:riemann}
$\zetaC:\Cstar\to\Cstar$ is the meromorphic continuation of $\zeta$ regarded as a map to
the Riemann sphere: $\zetaC(s)=\zeta(s)$ for $s\in\CC\setminus\{1\}$; $\zetaC(1)=\infty$
(the simple pole, Theorem~\ref{thm:riemann}); and $\zetaC(\infty)=1$, a \emph{conventional}
value at the adjoined domain point, chosen because $\zeta(s)\to1$ as $\re(s)\to+\infty$ (only
the $n=1$ term of $\sum_n n^{-s}$ survives). On the finite plane, $\zetaC$ is holomorphic into
$\Cstar$ on $\CC\setminus\{1\}$, its sole pole being $s=1$. It is \emph{not} continuous at the
domain point $\infty$, and no continuous extension exists there: $\zeta(-2n)=0$ for every
$n\ge1$ while $\zeta(\sigma)\to1$ as $\sigma\to+\infty$, so $\zeta$ has no limit at $\infty$.
Nothing downstream uses continuity at $\infty$: Definition~\ref{def:A-section} places its
conditions on the finite plane and on a slice half-space, and treats the value at $N$ as
assigned data. (Earlier versions of this definition asserted holomorphy on
$\Cstar\setminus\{1\}$; that assertion was false and was withdrawn on 2026-09-02.) The
functional equation, the conjugate symmetry
$\zetaC(\bar s)=\overline{\zetaC(s)}$, and Hardy's infinitude (Theorem~\ref{thm:hardy}) hold
verbatim for $\zetaC$, as they concern values on $\CC\setminus\{1\}$ where $\zetaC=\zeta$.
\end{definition}

\begin{lemma}[Compactification preserves the zero set]\label{lem:zero-Cstar}
\lean{zetaC_coe_eq_zero_iff, zetaC_zero_iff}
\uses{def:zeta-Cstar}
$Z(\zetaC)=Z(\zeta)$: the zeros of $\zetaC$ in $\Cstar$ are exactly the zeros of $\zeta$
in $\CC$; in particular the nontrivial zeros coincide.
\end{lemma}
\begin{proof}
On $\CC\setminus\{1\}$ the two functions agree, so they share all zeros there. The two
adjoined values lie outside the zero set: $\zetaC(1)=\infty$ is a pole and
$\zetaC(\infty)=1$. Hence compactification preserves the zero set exactly.
\end{proof}

%=====================================================================
\part{Slice-Preserving Theory, the Ring $\mathcal R$, and the Equivalence}
%=====================================================================

\section{Octonions and slices}

The octonions are an $8$-dimensional real division algebra, the largest of the normed division
algebras. This section fixes what we need from them: the algebra itself; Artin's theorem, that any
two octonions generate an associative subalgebra, so that powers and slicewise complex analysis
make sense; and the complex slices $\CC_v$, one for each imaginary direction $v\in S^6$, together
with their compactifications, the slice Riemann spheres. All the slices share the real axis and a
single point at infinity, the north pole $N$; that shared geometry is what the later construction
turns around.

The octonions $\OO$ have basis $\{1,e_1,\dots,e_7\}$ with each $e_i^2=-1$.

\begin{definition}\label{def:octonions}
\lean{Octonion, Octonion.normSq_mul, Octonion.alt_left, Octonion.alt_right}
The octonions are: (i) a \emph{division algebra}: every nonzero element has a
multiplicative inverse; (ii) a \emph{normed algebra}: $|ab|=|a||b|$; (iii)
\emph{non-associative}: $(ab)c\neq a(bc)$ in general; (iv) \emph{alternative}:
$(aa)b=a(ab)$ and $a(bb)=(ab)b$.
\end{definition}

An octonion $s\in\OO$ decomposes as $s=\sigma+w$ with $\sigma=\re(s)\in\RR$ and
$w=\im(s)\in\im(\OO)\cong\RR^7$; the unit imaginary octonions form $S^6\subset\im(\OO)$.

\begin{theorem}[Artin \cite{Schafer66}]\label{thm:artin}
In an alternative algebra, the subalgebra generated by any two elements is
associative.
\end{theorem}

\begin{corollary}\label{cor:powers}
\uses{thm:artin}
Powers $s^n$ are well-defined for any $s\in\OO$.
\end{corollary}

\begin{definition}[Complex slices and slice Riemann spheres]\label{def:slices}
\lean{Octonion.unitImaginarySphere, Octonion.sliceEmbed,
Octonion.sliceEmbed_mul, Octonion.sliceSphere}
\uses{def:octonions}
For any unit imaginary octonion $v\in S^6$, the \emph{complex slice} is
$\CC_v=\{a+bv:a,b\in\RR\}\subset\OO$. Since $v^2=-1$, this is a subalgebra isomorphic to
$\CC$ via the $\RR$-algebra isomorphism $\varphi_v:\CC\to\CC_v$, $i\mapsto v$. All slices
share the real axis: $\CC_v\cap\CC_w=\RR$ for $v\neq\pm w$. Its one-point compactification
is the \emph{slice Riemann sphere} $\CC_v^{*}:=\CC_v\cup\{\infty\}=S^2_v$, and $\varphi_v$
extends to an isomorphism of Riemann spheres $\varphi_v:\Cstar\xrightarrow{\sim}\CC_v^{*}$
with $\varphi_v(\infty)=\infty$ (GPV \cite{VS}, Prop.~4.1/Thm.~4.2). All slice spheres
share the single point at infinity, as well as the real axis.
\end{definition}

\section{Slice-preserving functions}
The class we use is the \emph{slice-preserving} slice-regular functions. We first recall slice
regularity \cite{GS07,GP11,CSS12,GSS13}, holomorphicity on each complex slice separately, as
the ambient notion, then isolate the slice-preserving subclass, on which classical complex methods
apply slicewise (by Artin's theorem each slice $\CC_v$ is an associative subalgebra where classical
complex analysis applies unchanged).

\begin{definition}[Slice regularity; {\cite[Def.~5.1.1]{CSS12}}]\label{def:slice-regular}
\uses{def:slices}
Let $\Omega\subseteq\OO$ be open. A function $f:\Omega\to\OO$ is \emph{slice regular}
if for each $v\in S^6$, $\varphi_v^{-1}\circ f\circ\varphi_v$ satisfies the
Cauchy--Riemann equations on $\varphi_v^{-1}(\Omega\cap\CC_v)\subseteq\CC$. A domain
$\Omega$ is \emph{axially symmetric} if $\sigma+\gamma v\in\Omega$ implies
$\sigma+\gamma w\in\Omega$ for all $w\in S^6$.
\end{definition}

\begin{theorem}[Representation Formula; {\cite[Thm.~5.1.7]{CSS12}}]\label{thm:rep-formula}
\uses{def:slice-regular}
Let $f$ be slice regular on an axially symmetric slice domain $\Omega$. Fix
$v_0\in S^6$ and write $f(\sigma+\gamma v_0)=\alpha(\sigma,\gamma)+v_0\cdot\beta(\sigma,\gamma)$
with $\alpha,\beta:\RR^2\to\RR$. Then for every $v\in S^6$,
$f(\sigma+\gamma v)=\alpha(\sigma,\gamma)+v\cdot\beta(\sigma,\gamma)$, where
$\alpha=\tfrac12[f(\sigma+\gamma v_0)+f(\sigma-\gamma v_0)]$ and
$\beta=\tfrac12 v_0^{-1}[f(\sigma-\gamma v_0)-f(\sigma+\gamma v_0)]$.
\end{theorem}

\begin{theorem}[Identity Theorem; {\cite[Cor.~5.1.9]{CSS12}}]\label{thm:identity}
\lean{stem_identity}
\uses{def:slice-regular}
Let $f$ be slice regular on an axially symmetric slice domain $\Omega$. If $f$
vanishes on a subset of a single slice $\Omega\cap\CC_{v_0}$ with an accumulation
point in $\Omega\cap\CC_{v_0}$, then $f\equiv0$ on all of $\Omega$.
\end{theorem}

\begin{theorem}[Extension Theorem; {\cite[Thm.~5.1.5]{CSS12}}]\label{thm:extension}
\uses{def:slice-regular}
Let $U\subseteq\CC$ be a domain symmetric under conjugation and $g:U\to\CC$
holomorphic. There is a unique slice regular $\mathrm{ext}(g):\Omega_U\to\OO$ on
$\Omega_U=\{\sigma+\gamma v:\sigma+i\gamma\in U,\ v\in S^6\}$ with
$\mathrm{ext}(g)|_{\CC_{v_0}}=\varphi_{v_0}\circ g\circ\varphi_{v_0}^{-1}$ for every
$v_0$.
\end{theorem}

\begin{definition}[Slice-preserving functions; \protect{\cite[Def.~2.7, Rem.~2.8]{AdF}; \cite{VS,SeriesExp,Wang}}]\label{def:slice-preserving}
\lean{IsIntrinsic, StemRing.isIntrinsic, StemRing.real_on_real}
\uses{def:slice-regular, thm:rep-formula}
A slice-regular function $f$ on an axially symmetric domain $\Omega$ is \emph{slice
preserving} if it maps every slice into itself: $f(\Omega\cap\CC_v)\subseteq\CC_v$ for all
$v\in S^6$. For octonions ($\#S^6>2$) the following are equivalent characterisations:
\begin{enumerate}[leftmargin=2.4em]
\item[(i)] $f$ is \emph{intrinsic}: $f(q^{c})=f(q)^{c}$ for all $q\in\Omega$, where $q^{c}$
  is octonionic conjugation (on a slice, $a+bv\mapsto a-bv$);
\item[(ii)] its representation-formula components are $\RR$-valued: writing
  $f(\sigma+\gamma v)=\alpha(\sigma,\gamma)+v\,\beta(\sigma,\gamma)$
  (Theorem~\ref{thm:rep-formula}), one has $\alpha,\beta:\RR^2\to\RR$, equivalently the
  inducing stem function $F=F_1+\iota F_2$ has $\RR$-valued components $F_1,F_2$;
\item[(iii)] $f$ preserves two distinct slices; equivalently $f(\Omega\cap\RR)\subseteq\RR$
  when $\Omega$ is a slice domain.
\end{enumerate}
Slice-preserving functions are the more restrictive subclass of slice-regular functions on
which classical complex methods apply slicewise. They are precisely the sections of
$\mathcal R$ (Definition~\ref{def:R}), and on them the regular product is the ordinary
pointwise product, so the product is commutative (Theorem~\ref{thm:wang}; Ghiloni--Perotti--
Stoppato \cite[Rem.~1.8]{SeriesExp}).
\end{definition}

\begin{remark}[Compactified extension]\label{rmk:slice-pres-compact}
Definition~\ref{def:slices} through Definition~\ref{def:slice-preserving}, and
Theorems~\ref{thm:rep-formula}--\ref{thm:extension}, are stated in the literature on
uncompactified domains $\Omega\subseteq\OO$ (and $U\subseteq\CC$). We use them on the
compactified $\Ostar=S^8$ and on the slice Riemann spheres $\CC_v^{*}=S^2_v$ (GPV's slice
Riemann manifolds \cite{VS}). The passage $\OO\rightsquigarrow\Ostar$,
$\CC_v\rightsquigarrow\CC_v^{*}$ adjoins only the shared point $\infty=N$; it is recorded as
a marked extension node (in the Lean development: cite the uncompactified statement and
extend through the one-point compactification \texttt{OnePoint}).
\end{remark}

\section{The octonionic zeta function \texorpdfstring{$\zetaO$}{zeta\_O}}

Slice-preserving theory already gives the tools to build the octonionic zeta function. On each slice
$\CC_v$ we read the compactified classical $\zeta$ through the isomorphism
$\varphi_v:\Cstar\to\CC_v^{*}$, and the representation formula glues these slice copies into a single
function on all of $\Ostar$. The one thing to check is that the two identifications of a slice with
$\CC$ agree, so that $\zetaO$ is well-defined; they do, because conjugating the input and reversing
the identification cancel. The pole of $\zeta$ at $s=1$ is carried to the north pole $N$.

\begin{definition}[Octonionic zeta]\label{def:zeta_O}
\lean{zetaO, zetaO_infty, zetaO_one, zetaO_sliceEmbed,
zetaO_sliceEmbed_of_im_neg, zetaO_ofReal}
\uses{def:zeta-Cstar, def:slices, thm:extension}
Here $N$ denotes the \emph{north pole} of $S^8$: the single point at infinity $\infty$
adjoined in the one-point compactification $\Ostar=\OO\cup\{\infty\}\cong S^8$ (inverse
stereographic projection; GPV \cite{VS}). It is one point with two names, so $\infty=N$;
it is the only point at infinity, shared by $\Ostar$ and by every slice Riemann sphere
$\CC_v^{*}=S^2_v$.

The \emph{octonionic zeta function} $\zetaO:\Ostar\to\Ostar$ is defined slicewise from the
compactified classical zeta (Definition~\ref{def:zeta-Cstar}) through the slice
identifications $\varphi_v:\Cstar\xrightarrow{\sim}\CC_v^{*}=S^2_v$ (the $\RR$-algebra
isomorphism $\CC\to\CC_v$ of Definition~\ref{def:slices}, extended to the Riemann spheres
by $\varphi_v(\infty)=N$; GPV \cite{VS}, Prop.~4.1/Thm.~4.2):
\begin{itemize}[leftmargin=2.2em]
\item[(i)] for $s\in\OO\setminus\RR$, with $w=\im(s)$, $v=w/|w|\in S^6$, $\gamma=|w|>0$,
  $\sigma=\re(s)$, so $s=\sigma+\gamma v$:\quad $\zetaO(s):=\varphi_v(\zetaC(\sigma+i\gamma))$;
\item[(ii)] for $s\in\RR\setminus\{1\}$:\quad $\zetaO(s):=\zetaC(s)=\zeta(s)\in\RR$;
\item[(iii)] for $s=1$ (the simple pole, Definition~\ref{def:zeta-Cstar}):\quad
  $\zetaO(1):=\infty=N$;
\item[(iv)] for $s=\infty=N$:\quad $\zetaO(\infty):=\zetaC(\infty)=1$.
\end{itemize}
\end{definition}

\begin{proposition}[Well-definedness]\label{prop:well-defined}
\lean{zetaO_sliceEmbed', riemannZeta_conj}
\uses{def:zeta_O, thm:riemann}
$\zetaO:\Ostar\to\Ostar$ is well-defined.
\end{proposition}
\begin{proof}
For $s\in\OO\setminus\RR$ the decomposition $s=\sigma+\gamma v$ with $\gamma>0$,
$v\in S^6$ is forced ($\gamma=|\im s|$, $v=\im s/|\im s|$), so the only ambiguity is the
identification of the slice $\CC_v=\CC_{-v}$ with $\Cstar$. That slice carries two such
identifications, $\varphi_v$ ($i\mapsto v$) and $\varphi_{-v}$ ($i\mapsto-v$), under
which $s$ reads as $\sigma+i\gamma$ resp.\ $\sigma-i\gamma$; we must check
$\varphi_v(\zetaC(\sigma+i\gamma))=\varphi_{-v}(\zetaC(\sigma-i\gamma))$. By conjugate
symmetry $\zetaC(\sigma-i\gamma)=\overline{\zetaC(\sigma+i\gamma)}$
(\cite[Ch.~2]{Titchmarsh86}; Definition~\ref{def:zeta-Cstar}), writing
$\zetaC(\sigma+i\gamma)=a+bi$ gives $\varphi_{-v}(a-bi)=a+(-b)(-v)=a+bv=\varphi_v(a+bi)$:
the two sign reversals, one from conjugating the input, one from the reversed
identification, cancel. The remaining cases carry no ambiguity: for real $s\neq1$,
$\zetaO(s)=\zeta(s)\in\RR\subseteq\CC_v^{*}$ for every $v$; at $s=1$ the value is the
single point $\infty=N$; and at $s=\infty=N$ the value is $\zetaC(\infty)=1\in\RR$, again
the same for every $v$ since each $\varphi_v$ fixes $\RR$ pointwise.
\end{proof}

\section{\texorpdfstring{$\mathcal R$}{R} is a commutative ring}\label{sec:R-ring}

The slice-preserving sections form a commutative ring $\mathcal R$, with $\zetaO$ among its elements.
Its multiplication is the regular ($\ast$) product, which descends from the Cayley--Dickson
construction of $\OO$ and from slice regularity; on slice-preserving functions it restricts on each
slice to the ordinary pointwise product (Wang), and is therefore commutative. The ring axioms, and
the integral-domain property, follow slicewise.

We work on $\Ostar$; write $\Omega_v:=\Ostar\cap\CC_v=\CC_v^{*}=S^2_v$ for the slice
Riemann sphere (Definition~\ref{def:slices}).

\begin{definition}[The ring of slice-preserving functions under the regular product]\label{def:R}
\lean{StemRing}
\uses{def:slice-preserving, def:slices}
$\mathcal R:=\mathcal O_{\Ostar}$ is the set of all slice-regular functions
$f:\Ostar\to\Ostar$ that are slice preserving (Definition~\ref{def:slice-preserving}),
$f(\Omega_v)\subseteq\CC_v^{*}$ for all $v\in S^6$, equipped with pointwise addition and
the regular ($\ast$) product.
\end{definition}

Commutativity of $\mathcal R$ rests on a single slice-preserving fact from the literature: on each
slice the regular product is ordinary multiplication.

\begin{theorem}[Regular product on slice-preserving functions; {\cite[Rem.~2.11]{Wang}}]\label{thm:wang}
\uses{thm:artin, def:R}
For slice-preserving $f,g$ (with $f(\Omega_v)\subseteq\CC_v$), the regular product
restricts on each slice to the pointwise product, $(f\ast g)|_{\Omega_v}(w)=f(w)\cdot
g(w)\in\CC_v$, and $f\ast g=g\ast f$, $(f\ast g)\ast h=f\ast(g\ast h)$ (via Artin,
Theorem~\ref{thm:artin}).
\end{theorem}

\begin{proposition}\label{prop:R-comm-ring}
\lean{StemRing}
\uses{thm:wang, thm:identity, def:R}
$(\mathcal R,+,\ast)$ is a commutative ring with identity $1\in\mathcal R$.
\end{proposition}
\begin{proof}
Closure under $+$ is pointwise. For $\ast$: fixing a slice $\CC_K$, the product
reduces to pointwise multiplication on $\CC_K$ by Theorem~\ref{thm:wang}, which is
commutative because $\CC_K$ is a commutative associative subalgebra; associativity
and distributivity are inherited slicewise, and a slice-regular function is
determined by its restriction to any one slice (Theorem~\ref{thm:identity}). The
constant $1$ is the identity.
\end{proof}

\begin{proposition}[$\mathcal R$ is an integral domain]\label{prop:R-domain}
\lean{StemRing.eq_zero_or_eq_zero_of_mul_eq_zero}
\uses{thm:wang, thm:identity}
If $f,g\in\mathcal R$ and $f\ast g=0$, then $f=0$ or $g=0$.
\end{proposition}
\begin{proof}
$(f\ast g)|_{\Omega_K}(w)=f(w)\cdot g(w)$ in $\CC_K\cong\CC$
(Theorem~\ref{thm:wang}). The restrictions $f|_{\Omega_K}$, $g|_{\Omega_K}$ are
holomorphic on the connected slice; by the classical identity principle the pointwise
product vanishes identically iff one factor does. By Theorem~\ref{thm:identity} the
corresponding global function vanishes.
\end{proof}

\section{The slice-preserving analytic toolkit on the compactified octonions}\label{sec:toolkit}

We collect the slice-preserving analytic facts used in Part~3 as cited statements: the exponential,
its logarithm manifold, its degenerate fibre, the slice/stem square and $I$-independent lift, and the
winding lift. Each is established in the literature on the uncompactified $\OO$ and used here on
$\Ostar=S^8$ and the slice Riemann spheres $\CC_v^{*}=S^2_v$ (the passage is recorded in
Remark~\ref{rmk:slice-pres-compact}; the compactified slice-manifold structure itself is GPV's \cite{VS}).

\begin{theorem}[The slice exponential; \protect{\cite[\S3]{GPVwind}; \cite[Prop.~5.1]{VS}; \cite[Rem.~2.23]{AdF}}]\label{thm:slice-exp}
\lean{Octonion.exp, Octonion.exp_sliceEmbed, Octonion.norm_exp,
Octonion.exp_ne_zero}
\uses{def:slice-preserving, cor:powers}
The exponential $\exp:\OO\to\OO\setminus\{0\}$, $\exp(q)=\sum_{k\ge0}q^{k}/k!$ (powers well defined by Corollary~\ref{cor:powers}), is a
slice-regular, slice-preserving entire function (Definition~\ref{def:slice-preserving}),
given on each slice by
\[
\exp(x+Iy)=e^{x}(\cos y+I\sin y),\qquad x,y\in\RR,\ I\in S^6 .
\]
Consequently $|\exp q|=e^{\re q}$ depends only on the real part, and $\exp$ is nowhere zero.
\end{theorem}

\begin{theorem}[The logarithm manifold $E^{+}_{\OO}$; \protect{\cite[Prop.~5.1, Rem.~5.2, Def.~5.3, Prop.~5.4, Def.~5.5]{VS}; \cite[\S3]{GPVwind}}]\label{thm:log-manifold}
\lean{Octonion.Eexp, Octonion.logManifold, Octonion.Llog,
Octonion.Llog_Eexp, Octonion.Eexp_Llog, Octonion.exp_Llog}
\uses{thm:slice-exp}
Let $T:\OO\to\OO\times\im\OO$, $T(x+Iy)=(\sinh x\cos y+I\sinh x\sin y,\ Iy)$, and
$\OO^{+}=\{q:\re q>0\}$. The \emph{logarithm manifold} is $E^{+}_{\OO}:=T(\OO^{+})$, the
semi-helicoidal hypercomplex $8$-manifold; equivalently it is the image of the exponential
graph, parametrised by the \emph{$E^{+}_{\OO}$-exponential}
\[
E:\OO\longrightarrow E^{+}_{\OO}\subset\OO\times\im\OO,\qquad
E(x+Iy)=(\exp(x+Iy),\,Iy)=(e^{x}\cos y+Ie^{x}\sin y,\ Iy),
\]
an immersion and a diffeomorphism, with inverse the \emph{$E^{+}_{\OO}$-logarithm}
$L:E^{+}_{\OO}\to\OO$, $L(q,p)=\log|q|+p$ (here $p=\mathrm{Arg}(q)$). Writing $\pi$ for the
projection to the first factor, $\pi\circ E=\exp$ (Theorem~\ref{thm:slice-exp}). The manifold
$E^{+}_{\OO}$ is an adapted blow-up of $\OO$ along the real axis. Its projection
$\pi:E^{+}_{\OO}\to\OO\setminus\{0\}$ has the degenerate real fibres responsible for the
failure of covering and openness. The blow-up unwraps the multivalued logarithm into the
single-valued $L$, and a branch of the
hypercomplex logarithm is $\log_{\OO}=L\circ(\pi|_{\Omega})^{-1}$ on any path-connected
$\Omega\subset E^{+}_{\OO}$ on which $\pi|_{\Omega}$ is injective.
\end{theorem}

\begin{lemma}[The degenerate set of the exponential; its fibre over a real value]\label{lem:exp-degenerate}
\lean{Octonion.exp_fibre_neg_real, Octonion.exp_eq_neg_ofReal_iff,
Octonion.logManifold_fibre_neg_real, Octonion.degenerate_level_readout}
\uses{thm:slice-exp, thm:log-manifold}
\emph{Sourced.} With $\pi$ the first-factor projection, $\pi\circ E=\exp$
\textup{(\cite[Rem.~5.2(a)]{VS}, the commuting triangle)}. Its spherical
\emph{degenerate set} makes the projection an adapted blow-up with failures of covering and
openness \textup{(\cite[Rem.~5.2(b)]{VS})}. The direction $I(q)$ becomes singular at every
point of $\RR$ \textup{(\cite[Rem.~2.1]{GPVwind})}.

\emph{Proved from the slice form.} For $r>0$,
\[
\exp^{-1}(-r)=\bigl\{\,\log r+I(2k{+}1)\pi:\ I\in S^6,\ k\in\ZZ\,\bigr\}:
\]
by Theorem~\ref{thm:slice-exp}, $\exp(x+Iy)=e^{x}\cos y+Ie^{x}\sin y=-r$ forces
$e^{x}\sin y=0$, hence $y\in\pi\ZZ$; negativity of the value forces $y=(2k{+}1)\pi$ and
$x=\log r$; the direction $I$ is unconstrained. The fibre is thus indexed by the single real
\emph{level} $\log r=\log|{-}r|$ and, within it, by the odd winding indices $2k{+}1$ carried as
band data (the winding lift, Theorem~\ref{thm:winding-lift}; \cite[Cor.~5.21]{GPVwind}).
The fibre formula itself appears, stated without proof as Preface motivation, in \cite{VS}
(p.~972); the slice-form derivation above is the load-bearing statement.
\end{lemma}

\begin{definition}[Slice restriction and the $I$-independent lift; \protect{\cite[Rem.~2.11]{Wang}; \cite[\S3.3,~\S3.7]{BisiWinkelmann}; \cite[Def.~2.1, Prop.~2.2]{SeriesExp}}]\label{def:slice-squares}
\lean{IsIntrinsic, ASection.realize, ASection.realize_mem_sliceSphere,
ASection.realize_equivariant}
A slice-preserving section $A$ carries each slice Riemann sphere into itself,
\[
  A(\CC_v^{*})\subseteq\CC_v^{*}=S^2_v\qquad(v\in S^6),
\]
so its restriction to a slice is a holomorphic self-map of that slice sphere,
\[
  A_v:=\varphi_v^{-1}\circ A\circ\varphi_v:\ \Cstar\longrightarrow\Cstar,
\]
through the identification $\varphi_v:\Cstar\xrightarrow{\sim}\CC_v^{*}$ of
Definition~\ref{def:slices}. The \emph{$I$-independent lift} is the fact that this holomorphic stem is
the \emph{same} for every $v$: there is one holomorphic $A_\circ:\Cstar\to\Cstar$ with
\[
  A|_{\CC_v^{*}}=\varphi_v\circ A_\circ\circ\varphi_v^{-1}\qquad\text{for all }v\in S^6
\]
(Wang \cite[Rem.~2.11]{Wang}; Bisi--Winkelmann \cite[\S3.7]{BisiWinkelmann}). Equivalently, $A$ is
\emph{intrinsic}, $A(q^c)=A(q)^c$, and is induced ($A=\mathcal I(F)$) by a single stem function
$F=F_1+\iota F_2:D\to\OO_{\CC}$ on the symmetric $D\subset\CC$ whose components $F_1,F_2$ are
$\RR$-valued (the stem-function formalism, Ghiloni--Perotti \cite[Def.~2.1, Prop.~2.2]{SeriesExp}):
\[
  A(x+Iy)=F_1(x+iy)+I\,F_2(x+iy),\qquad F_1,F_2\ \RR\text{-valued}.
\]
It is this single $I$-independent slice stem that makes a slice-preserving section
$\Gtwo$-equivariant, hence blind to the $\Gtwo$-orbit direction
(Definition~\ref{def:section-map}).
\end{definition}

\begin{theorem}[The winding lift; \protect{\cite[Def.~3.4, Def.~4.1, Prop.~4.2, Def.~5.11]{GPVwind}}]\label{thm:winding-lift}
\lean{exists_log_continuation, winding_lift_unique,
Octonion.lift_iff_continuation}
\uses{thm:log-manifold}
Let $\gamma:[a,b]\to\OO\setminus\{0\}$ be a path \textup{(}on each slice $\CC_v$ this is the
classical complex continuation\textup{)}. A \emph{continuation of the hypercomplex logarithm along
$\gamma$} is a path $\tilde\gamma:[a,b]\to\OO$ with $\exp\circ\tilde\gamma=\gamma$, and a \emph{lift
of $\gamma$} to the logarithm manifold $E^{+}_{\OO}$ \textup{(}Theorem~\ref{thm:log-manifold}\textup{)}
is a path $\Gamma:[a,b]\to E^{+}_{\OO}$ with $\mathrm{pr}_1\circ\Gamma=\gamma$:
\[
\begin{tikzcd}[row sep=1.6em, column sep=2.6em]
& \OO\\
{[a,b]} \arrow[ur,"\tilde\gamma"] \arrow[r,"\gamma"'] & \OO\setminus\{0\} \arrow[u,"\exp"']
\end{tikzcd}
\qquad\qquad
\begin{tikzcd}[row sep=1.6em, column sep=2.6em]
& E^{+}_{\OO} \arrow[d,"\mathrm{pr}_1"]\\
{[a,b]} \arrow[ur,"\Gamma"] \arrow[r,"\gamma"'] & \OO\setminus\{0\}
\end{tikzcd}
\]
The two data are \emph{equivalent}: a lift $\Gamma$ exists if and only if a continuation
$\tilde\gamma$ exists, and they correspond by
\[
\tilde\gamma=L\circ\Gamma,\qquad \Gamma=(\exp\circ\tilde\gamma,\ \im\tilde\gamma),
\]
where $L$ is the $E^{+}_{\OO}$-logarithm of Theorem~\ref{thm:log-manifold}
\textup{(}\cite[Prop.~4.2]{GPVwind}\textup{)}. For a \emph{loop} $\gamma:S^1\to\OO\setminus\{0\}$ the
lift is taken on the universal cover $\exp:i\RR\to S^1$, as a continuous $\Gamma:i\RR\to E^{+}_{\OO}$
with $\mathrm{pr}_1\circ\Gamma=\gamma\circ\exp$ \textup{(}\cite[Def.~5.11]{GPVwind}\textup{)}:
\[
\begin{tikzcd}[row sep=1.6em, column sep=2.6em]
i\RR \arrow[r,"\Gamma"] \arrow[d,"\exp"'] & E^{+}_{\OO} \arrow[d,"\mathrm{pr}_1"]\\
S^1 \arrow[r,"\gamma"'] & \OO\setminus\{0\}
\end{tikzcd}
\]
\end{theorem}

\begin{proposition}[Tameness, existence, and the winding number; \protect{\cite[Def.~4.20, Cor.~5.21, Cor.~5.22]{GPVwind}}]\label{prop:winding-signature}
\lean{winding_lift_unique, winding_loop_defect}
\uses{thm:winding-lift}
The lift of Theorem~\ref{thm:winding-lift} is governed by the companion structure of $\gamma$.
\begin{enumerate}[label=\textup{(\roman*)},leftmargin=2.6em]
\item \emph{Uniqueness (tameness).} $\gamma$ is \emph{tame} when it admits a \emph{unique} companion
  imaginary-unit field; its lift is then unique \textup{(}\cite[Def.~4.20]{GPVwind}\textup{)}.
\item \emph{Existence and closure.} A lift of a loop $\gamma$ exists if and only if the signature
  satisfies $\sigma(\gamma|_{[\xi_l,\xi_{l+1}]})\in\{0,-1\}$ on each obstruction interval; and
  \emph{if it exists, the lift is itself a loop} \textup{(}\cite[Cor.~5.13]{GPVwind}; pinned
  against the arXiv text; the earlier citation to Cor.~5.22 is superseded\textup{)}.
\item \emph{Winding number.} For an untwisted tame loop with even circular signature $\sigma^c$, the
  winding number is $\omega=|\sigma^c|/2$ \textup{(}\cite[Cor.~5.21]{GPVwind}\textup{)}.
\end{enumerate}
\end{proposition}

\section{Slice preservation and symmetries}

With the ring in hand, we check that $\zetaO$ actually lives in it, and record the symmetry that will
organize its zeros. Slice preservation says $\zetaO$ keeps each slice inside itself; together with
slice regularity this makes it a section of $\mathcal R$. The symmetry is the action of the
automorphism group $\Gtwo=\Aut(\OO)$, which fixes the real axis and rotates the imaginary directions
$S^6$, and $\zetaO$ is equivariant under it. This forces every non-real zero to fill a whole
$6$-sphere, the geometry the next section reads off.

\begin{theorem}[Slice preservation]\label{thm:slice-pres}
\lean{zetaO_mem_sliceSphere}
\uses{def:zeta_O, def:slice-preserving}
$\zetaO$ is slice preserving (Definition~\ref{def:slice-preserving}):
$\zetaO(\CC_v^{*})\subseteq\CC_v^{*}$ for every $v\in S^6$.
\end{theorem}
\begin{proof}
Fix $v\in S^6$ and $s\in\CC_v^{*}$. If $s\in\CC_v\setminus\RR$ then $s=\sigma\pm\gamma v$
and $\zetaO(s)=\varphi_v(\zetaC(\sigma\pm i\gamma))\in\varphi_v(\Cstar)=\CC_v^{*}$
(Definition~\ref{def:zeta_O}(i)). If $s\in\RR\setminus\{1\}$ then
$\zetaO(s)=\zeta(s)\in\RR\subseteq\CC_v^{*}$; while $\zetaO(1)=\infty\in\CC_v^{*}$ and
$\zetaO(\infty)=1\in\CC_v^{*}$. In every case the value remains in $\CC_v^{*}$.
\end{proof}

\begin{theorem}[$\zetaO$ is a section of $\mathcal R$]\label{thm:zeta-in-R}
\lean{zetaSection}
\uses{thm:slice-pres, thm:extension, def:R, prop:well-defined}
$\zetaO\in\mathcal R$.
\end{theorem}
\begin{proof}
On each slice the restriction is the compactified classical zeta read through the
isomorphism $\varphi_v$, namely $\zetaO|_{\CC_v^{*}}=\varphi_v\circ\zetaC\circ\varphi_v^{-1}$,
holomorphic into $\CC_v^{*}$ away from the single pole.
Equivalently $\zetaO=\mathrm{ext}(\zeta)$ is the slice-regular extension of the
holomorphic, conjugation-symmetric $\zeta$ on $\CC\setminus\{1\}$
(Theorem~\ref{thm:extension}, \cite[Thm.~5.1.5]{CSS12}); hence $\zetaO$ is slice regular
on $\Ostar\setminus\{1,\infty\}$ (slice regularity is slicewise holomorphy,
Definition~\ref{def:slice-regular}). By Theorem~\ref{thm:slice-pres} it is slice
preserving, $\zetaO(\CC_v^{*})\subseteq\CC_v^{*}$. A slice-regular, slice-preserving
function $\Ostar\to\Ostar$ is by definition a section of $\mathcal R=\mathcal O_{\Ostar}$
(Definition~\ref{def:R}); the regular product on such sections restricts on each slice to
the ordinary product (Theorem~\ref{thm:wang}, Wang \cite[Rem.~2.11]{Wang}). Its zeros
\textup{(}Theorem~\ref{thm:zero-spheres}\textup{)} make $\zetaO$ a nonunit.
\end{proof}

\begin{definition}[The automorphism group $\Gtwo$]\label{def:G2}
\lean{G2, G2.smul_one, G2.smul_ofReal}
$\Gtwo=\Aut(\OO)$ is the $14$-dimensional compact, connected, simply connected Lie group of
$\RR$-algebra automorphisms of $\OO$. Each $g\in\Gtwo$ fixes $1$, hence the real axis pointwise, and
acts on $\im\OO\cong\RR^7$ as an element of $SO(7)$; thus $g\cdot(\sigma+\gamma v)=\sigma+\gamma\,g(v)$.
\end{definition}

\begin{theorem}[$\Gtwo$-equivariance]\label{thm:G2-equiv}
\lean{zetaO_equivariant}
\uses{def:G2, def:zeta_O}
$\zetaO(g\cdot s)=g\cdot\zetaO(s)$ for all $g\in\Gtwo$, $s\in\OO$.
\end{theorem}
\begin{proof}
For $s\in\RR$, $g\cdot s=s$ and $\zetaO(s)\in\RR$ is fixed, so both sides equal $\zetaO(s)$. For a
non-real $s=\sigma+\gamma v$ ($v\in S^6$, $\gamma>0$), $g$ carries the slice $\CC_v$ to the slice
$\CC_{g(v)}$ by the isometric relabelling $\varphi_{g(v)}=g\circ\varphi_v$
(Definition~\ref{def:G2}). Since $\zetaO$ is defined slicewise by the same $\zetaC$ read through
$\varphi_v$ (Definition~\ref{def:zeta_O}(i)),
\[
  \zetaO(g\cdot s)=\varphi_{g(v)}\bigl(\zetaC(\sigma+i\gamma)\bigr)
                 =g\bigl(\varphi_{v}(\zetaC(\sigma+i\gamma))\bigr)=g\cdot\zetaO(s).
\]
\end{proof}

\begin{theorem}[\protect{\cite{Baez02}}]\label{thm:G2-S6}
\lean{G2.exists_smul_eq_of_mem_unitImaginarySphere}
$\Gtwo$ acts transitively on $S^6$ with stabilizer $\mathrm{SU}(3)$:
$\mathrm{SU}(3)\hookrightarrow\Gtwo\twoheadrightarrow S^6$.
\end{theorem}

\begin{remark}[The $\Gtwo$-action extends to $\Ostar=S^8$]\label{rmk:G2-compact}
\uses{def:G2, def:carrier}
The action above is stated on $\OO$, but $\mathcal H_1$ (Part~3) runs on the compactified
$\Ostar=S^8$. Each $g\in\Gtwo$ fixes the real axis pointwise (an automorphism fixes $1$, hence all
of $\RR$) and acts on $\im\OO=\RR^7$ as an element of $SO(7)$; in particular $g$ is a linear
isometry of $\RR^8=\OO$, hence proper, and therefore extends uniquely to a homeomorphism of the
one-point compactification $\Ostar=\OO\cup\{\infty\}=S^8$ fixing the point at infinity $\infty=N$ (the one-point
compactification is functorial on proper maps; in Lean, \texttt{OnePoint}). Thus
$\Gtwo\curvearrowright\Ostar$ fixes the compactified real axis $\RR\cup\{\infty\}$, a great circle
through $N$, pointwise, and acts on each direction-sphere $S^6$ exactly as in
Theorem~\ref{thm:G2-S6}. This is the action on which $\mathcal H_1$ is built, with $\mathcal R$ its
ring of invariants over $\Ostar$.
\end{remark}

\section{Zero geometry and the equivalence}

This section is the bridge the whole paper is built to cross. It shows that the zeros of $\zetaO$ are
exactly the classical nontrivial zeros, now spread out: each nontrivial zero $\rho$ of $\zeta$
becomes a whole $6$-sphere of zeros of $\zetaO$, its residue-$\CC$ zero-sphere, sitting at real part
$\re\rho$. The Riemann Hypothesis, that every $\rho$ has real part $\tfrac12$, becomes the statement
that all these $6$-spheres share the same real part, that they are concentric. From here on we work
with the spheres.

\begin{theorem}[Zero Equivalence Theorem]\label{thm:zero-equivalence}
\lean{sliceEmbed_eq_zero_iff, zetaO_zero_iff}
\uses{def:zeta_O, thm:slice-pres, lem:zero-Cstar}
Let $\rho=\sigma+i\gamma\in\Cstar$. For every $v\in S^6$,
\[
\zetaO(\sigma+\gamma v)=0\quad\Longleftrightarrow\quad\zetaC(\rho)=0 .
\]
Equivalently: a non-real $s\in\Ostar$ is a zero of $\zetaO$ if and only if its slice
coordinate $\varphi_v^{-1}(s)$ is a zero of $\zetaC$; by Lemma~\ref{lem:zero-Cstar} these
are exactly the nontrivial zeros of the classical $\zeta$.
\end{theorem}
\begin{proof}
By Definition~\ref{def:zeta_O}(i), $\zetaO(\sigma+\gamma v)=\varphi_v(\zetaC(\sigma+i\gamma))$.
The slice identification $\varphi_v:\Cstar\to\CC_v^{*}=S^2_v$ is an isomorphism of slice
Riemann spheres (the $\RR$-algebra isomorphism $\CC\to\CC_v$ of
Definition~\ref{def:slices}, extended by $\varphi_v(\infty)=N$; GPV \cite{VS}). An
isomorphism is injective and sends $0$ to $0$, so $\varphi_v(z)=0\iff z=0$; with
$z=\zetaC(\rho)$ this gives both implications. (Slice preservation,
Theorem~\ref{thm:slice-pres}, is what places the value in $\CC_v^{*}$ so that
$\varphi_v^{-1}$ applies on the non-real side.) The final clause is
Lemma~\ref{lem:zero-Cstar}.
\end{proof}

\begin{theorem}[Nontrivial zeros are infinitely many disjoint $6$-spheres]\label{thm:zero-spheres}
\lean{zeroSphere, zeroSphere_eq_orbit, zeroSphere_zeros,
zeroSphere_disjoint, zeroSpheres_infinite}
\uses{thm:zero-equivalence, thm:G2-S6, thm:G2-equiv, cor:hadamard-infinitude}
For a nontrivial zero $\rho=\sigma+i\gamma$ of $\zetaC$ (so $\gamma>0$) set
\[
S_\rho=\{\sigma+\gamma v:v\in S^6\}=\sigma+\gamma\,S^6\subset\OO .
\]
Then:
\begin{enumerate}[leftmargin=2.4em]
\item[(i)] $S_\rho$ \emph{is} a $6$-sphere, the round sphere of radius $\gamma$ centred
  at the real point $\sigma$, equivalently the $\Gtwo$-orbit of any one of its points
  $\sigma+\gamma v_0$ (Theorem~\ref{thm:G2-equiv}; $\Gtwo$ acts transitively on $S^6$ with
  stabiliser $\mathrm{SU}(3)$, Theorem~\ref{thm:G2-S6}, Baez \cite{Baez02}).
\item[(ii)] every point of $S_\rho$ is a zero of $\zetaO$, and conjugate zeros give the
  same sphere, $S_\rho=S_{\bar\rho}$;
\item[(iii)] the nontrivial-zero set is the disjoint union
  $Z^{\mathrm{nt}}_{\Ostar}=\bigsqcup_{[\rho]}S_\rho$ over conjugate pairs
  $[\rho]=\{\rho,\bar\rho\}$, with distinct pairs giving disjoint spheres;
\item[(iv)] there are infinitely many such spheres (Corollary~\ref{cor:hadamard-infinitude};
  also Hardy \cite{Hardy14}, Theorem~\ref{thm:hardy}).
\end{enumerate}
\end{theorem}
\begin{proof}
(i) The map $v\mapsto\sigma+\gamma v$ is the affine map $x\mapsto\sigma+\gamma x$ of
$\im\OO\cong\RR^7$ restricted to $S^6$; with $\gamma>0$ its image is the round $6$-sphere
of radius $\gamma$ centred at $\sigma$. By $\Gtwo$-equivariance
(Theorem~\ref{thm:G2-equiv}) and transitivity of $\Gtwo$ on $S^6$
(Theorem~\ref{thm:G2-S6}), $\Gtwo\cdot(\sigma+\gamma v_0)=\{\sigma+\gamma\,g(v_0):g\in\Gtwo\}=S_\rho$.
So $S_\rho$ \emph{is} $S^6$, scaled by $\gamma$, translated to $\sigma$, realized as
a single $\Gtwo$-orbit; no embedding or diffeomorphism is invoked.

(ii) For any $v\in S^6$, $\zetaO(\sigma+\gamma v)=\varphi_v(\zetaC(\rho))=\varphi_v(0)=0$
by the Zero Equivalence Theorem~\ref{thm:zero-equivalence}. And
$S_{\bar\rho}=\{\sigma+\gamma(-v):v\in S^6\}=S_\rho$, since $v\mapsto-v$ is a bijection of
$S^6$.

(iii) For ``$\subseteq$'': a non-real zero $s=\sigma+\gamma v$ of $\zetaO$ forces
$\zetaC(\sigma+i\gamma)=0$ by Theorem~\ref{thm:zero-equivalence}, so $s\in S_\rho$ with
$\rho=\sigma+i\gamma$ nontrivial (Lemma~\ref{lem:zero-Cstar}); ``$\supseteq$'' is (ii).
Disjointness: if $\sigma+\gamma v=\sigma'+\gamma'w$, uniqueness of the real/imaginary
parts gives $\sigma=\sigma'$ and $\gamma v=\gamma'w$, whence $\gamma=\gamma'$ (taking
norms, $|v|=|w|=1$) and $\{\rho',\bar\rho'\}=\{\rho,\bar\rho\}$; so two spheres coincide or
are disjoint.

(iv) There are infinitely many spheres because $\zeta$ has infinitely many nontrivial zeros
(Corollary~\ref{cor:hadamard-infinitude}; also Hardy, Theorem~\ref{thm:hardy}).
\end{proof}

% [Conjugate pairs, Disjointness, and Zero-Set Decomposition are folded into
%  Theorem~\ref{thm:zero-spheres} above (items (ii)--(iv) and its proof).]

\begin{theorem}[The equivalence: concentricity $\Leftrightarrow$ RH]\label{thm:rh-equiv}
\lean{riemannHypothesis_iff_concentric}
\uses{thm:zero-spheres, thm:riemann}
The following are equivalent:
\begin{enumerate}[leftmargin=2.4em]
\item[(a)] the Riemann Hypothesis;
\item[(b)] the nontrivial-zero $6$-spheres $\{S_\rho\}$ of $\zetaO$
  (Theorem~\ref{thm:zero-spheres}) are \emph{concentric}: they share a single common
  centre, necessarily a point of the real axis.
\end{enumerate}
When they hold, the common centre is $\tfrac12$.
\end{theorem}
\begin{proof}
Each $S_\rho$ is centred at the real point $\sigma_\rho$, where
$\rho=\sigma_\rho+i\gamma_\rho$ (Theorem~\ref{thm:zero-spheres}(i)).

$(a)\Rightarrow(b)$: under RH every $\sigma_\rho=\tfrac12$, so all spheres share the
centre $\tfrac12$.

$(b)\Rightarrow(a)$: suppose the spheres share a common centre $c\in\RR$. By the functional
equation (Theorem~\ref{thm:riemann}), if $\rho$ is a nontrivial zero then so is
$1-\bar\rho=(1-\sigma_\rho)+i\gamma_\rho$, whose sphere $S_{1-\bar\rho}$ is centred at
$1-\sigma_\rho$. Concentricity forces $\sigma_\rho=c=1-\sigma_\rho$, hence
$\sigma_\rho=\tfrac12$; as $\rho$ was arbitrary this is RH, and the common centre is $\tfrac12$.
\end{proof}

%=====================================================================
\part{The Categorical Construction and the Concentricity Question}\label{part:stack}
%=====================================================================

Part 2 revealed that $\zeta$ carries a special geometry: spread onto the octonions, its zeros are
concentric exactly when the Riemann Hypothesis holds. That raises the question this part studies.
What general properties of a section of the ring $\mathcal R$ of slice-preserving functions force
its residue-$\CC$ zeros to be concentric? Following that question leads to four conditions,
\textup{C1--C4}, and to the whole class of functions satisfying them, which we call the A-sections.
The answer this part now records is that \textup{C1--C4} do \emph{not} suffice: the class
contains $(q^{2}+1)\,\zetaO$, whose zero-spheres are centred at $0$ and at the real parts of the
zeta zeros (Proposition~\ref{prop:countermodel}). What the categorical construction proves for
every A-section is that the residue-$\CC$ zero-sphere states form one connected total with a
singleton component colimit (Theorem~\ref{thm:residue-total-singleton}); what it leaves open is
whether the real reads of those states agree, and that is the concentricity question itself
(Proposition~\ref{prop:read-equivalence}). For the octonionic zeta that question is the Riemann
Hypothesis (Theorem~\ref{thm:zeta-concentricity-rh}).

From the section's own data we read one distinguished exponential M\"obius transformation of the
Riemann sphere, presented at its two fixed points by the Euler product and the Weierstrass product;
orbit--stabilizer carries that single transformation across the projective base. The resulting map
to $\SphereWorld$ is its geometric projection, while the same action on normalized A-value states
assembles the genuine functor $\mathcal A_A:\mathcal B\to\mathbf{Grpd}$. Its Grothendieck
construction $\mathcal T_A=\int_{\mathcal B}\mathcal A_A$ gathers every state and every
value-transport into one total object. The zero-sphere system is the inverse-image subdiagram
$\mathcal R_A$, included by $\iota_A:\mathcal R_A\Rightarrow\mathcal A_A$; the relative
orbit--stabilizer comparison makes $\int_{\mathcal B}\mathcal R_A$ transitive and hence connected.
The component colimit therefore collapses to a singleton. Each certified residue representative
carries a real read, the real coordinate of its zero-sphere; whether all those reads coincide is
not decided by the singleton, and Section~\ref{sec:detector} says exactly what is and is not
established.

\section{The geometric worlds and the categorical construction}\label{sec:tower-stack}

This section assembles the parts the detector runs on. We begin with the compactified octonions
$\Ostar=S^8$ and their north pole $N$, the point at infinity where a section sends its pole. We then
describe what a slice-preserving section does on each slice, one holomorphic stem repeated in every
direction, which is why its non-real zeros fill whole $6$-spheres. We set up the two geometric
viewpoints that carry this data, the octonionic world $\mathcal H_1$ and the slice-value world
$\mathcal S_2$. And we close with the projective base $\mathcal B$, the A-section functor over it,
and its Grothendieck total $\mathcal T_A$, the object the next section reads.

\subsection*{The compactified octonions and the point at infinity}

\begin{definition}[The compactified octonions $\OO^{*}=S^8$]\label{def:carrier}
\lean{Octonion, onePointEquivSphereOfFinrankEq}
Let $\OO^{*}=\OO\cup\{\infty\}=S^8$ be the Alexandroff one-point compactification of $\OO=\RR^8$.
Concretely (GPV \cite{VS}, Prop.~4.1 and Thm.~4.2, for $m=\dim_{\RR}K\in\{4,8\}$): the inverse
stereographic projections from the north and south poles,
\[
f(x+Iy)=\Big(\tfrac{2(x+Iy)}{1+x^2+y^2},\ \tfrac{-1+x^2+y^2}{1+x^2+y^2}\Big),\qquad
h(x+Iy)=\Big(\tfrac{2(x-Iy)}{1+x^2+y^2},\ \tfrac{1-x^2-y^2}{1+x^2+y^2}\Big),
\]
give a conformal atlas $\{(K,f),(K,h)\}$ endowing $S^m$ with the structure of a slice
(quaternionic/octonionic) Riemann manifold, with transition map $h^{-1}\circ f(q)=\bar q/q^2=1/q$,
itself slice regular and slice preserving. The adjoined point is the north pole $N$; this is the
point at infinity. In Lean the underlying compactification is
\texttt{Mathlib.\-Topology.\-Compactification.\-OnePoint}, $\OO^{*}=\texttt{OnePoint }\OO$ with
$\infty=\texttt{OnePoint.infty}$, and the homeomorphism $\OO^{*}\cong S^8$ is
\texttt{onePoint\-Equiv\-Sphere\-Of\-Finrank\-Eq} applied to $\dim_{\RR}\OO=8$
(\texttt{Mathlib.\-Topology.\-Compactification.\-OnePoint.\-Sphere}). The compactification is what makes
the value of a section at its real pole a well-defined point of $\OO^{*}$, the point at
infinity $\infty=N$. The two views $\mathcal H_1$ and $\mathcal S_2$ are built on
these compactified values (Definition~\ref{def:two-worlds}); the distinct projective base
$\mathcal B$ and direction-indexed sphere world used in the construction below are introduced in
Definition~\ref{def:base}.
\end{definition}

There is one geometric north pole $N=\infty$, seen from two categorical points of view. In
$\mathcal H_1$ this point of $\OO^{*}$ is a $\Gtwo$-fixed object: its own connected component,
with $\Aut(N)=\Gtwo$, carrying no universal property. In the slice-value world
$\mathcal S_2$ the same geometric point is represented by the object
\[
  \nfr\;:=\;\infty=N\ \text{regarded as an object of }\mathcal S_2,
\]
the common point at infinity of the slice Riemann spheres, fixed by band and direction alike.
The two representatives belong to \emph{different} categorical viewpoints and are related by the section
through \textup{C1}: the pole of $A$ has value $\infty$, so the section carries the pole point
to $\nfr$. The residue field of a closed zero sorts the picture:
residue-$\RR$ (the $\Gtwo$-fixed locus $\RR\cup\{N\}$, classically trivial) versus residue-$\CC$
(the $6$-sphere orbits, classically nontrivial; the dictionary is Part~2 and
Theorem~\ref{thm:zero-spheres}).
The later total construction uses this one shared compactified witness $N$, the same for every
zero index.

\subsection*{What a slice-preserving section does}

\begin{definition}[The section's slice data and equivariance]\label{def:section-map}
\lean{ASection.realize, ASection.realize_mem_sliceSphere,
ASection.realize_equivariant, ASection.sliceCoord_smul_invariant}
\uses{def:slice-squares, def:slice-preserving, thm:rep-formula}
Let $A\in\mathcal R$. By the $I$-independent lift (Definition~\ref{def:slice-squares}; Wang
\cite[Rem.~2.11]{Wang}, Bisi--Winkelmann \cite[\S3.7]{BisiWinkelmann}) there is one holomorphic
stem with $\RR$-valued components,
\[
  A(x+Iy)=F_1(x+iy)+I\,F_2(x+iy),\qquad F_1,F_2\ \RR\text{-valued},
\]
the same for every $I\in S^6$. Three consequences used below:
\begin{enumerate}[label=\textup{(\roman*)},leftmargin=2.6em]
\item \emph{Slice preservation of values.} $A(\CC_I^{*})\subseteq\CC_I^{*}$ for every $I$: the
  value at a point lies on that point's own slice sphere (Definition~\ref{def:R}).
\item \emph{$\Gtwo$-equivariance.} For $g\in\Gtwo$ and non-real $q=x+Iy$,
  \[
    A(g\cdot q)=F_1+g(I)\,F_2=g\cdot\bigl(F_1+I\,F_2\bigr)=g\cdot A(q),
  \]
  since $g\cdot(x+Iy)=x+g(I)y$ (Definition~\ref{def:G2}) and $F_1,F_2$ are real; for
  $q\in\RR\cup\{N\}$ both sides are $\Gtwo$-fixed, because $A(\RR)\subseteq\RR$ and
  $A(N)\in\bigcap_I\CC_I^{*}=\RR\cup\{N\}$ by \textup{(i)}. (For $\zetaO$ this is
  Theorem~\ref{thm:G2-equiv}; the computation above is the same one, run for an arbitrary
  intrinsic section.)
\item \emph{Blindness to the sphere direction.} On a $\Gtwo$-orbit $S_{(\sigma,\gamma)}$ the stem
  coordinates $(F_1,F_2)(\sigma+i\gamma)$ are constant; in particular if $A$ vanishes at one
  point of the orbit it vanishes on the whole orbit (Lemma~\ref{lem:residue-spheres}).
\end{enumerate}
A slice-preserving section carries one holomorphic stem per slice, and the direction $S^6$ is
traversed only by the $\Gtwo$-morphisms already present in both worlds.
\end{definition}

Condition \textup{C4}, infinitely many residue-$\CC$ zeros, matches the class to its
infinite structure (the infinite Euler family of \textup{C2}, the infinite Weierstrass divisor of
\textup{C3}) and gives the collapse its force: the degenerate fibre of the transport is an
infinite family, so that whether it lies in one connected component of $\mathcal T_A$
(Definition~\ref{def:base}) is an infinite alignment. Theorem~\ref{thm:residue-total-singleton}
decides the component question affirmatively for the C-residue total; the alignment of the real
reads is the separate question of Proposition~\ref{prop:read-equivalence}.

\subsection*{Two geometric viewpoints}

\begin{definition}[The octonionic and slice-value worlds]\label{def:two-worlds}
\lean{H1, SphereWorld}
\uses{def:carrier, def:G2, thm:G2-S6, rmk:G2-compact, def:slices}
Recall the \emph{translation groupoid} of an action $G\curvearrowright X$: its objects are the
points of $X$ and its morphisms $x\to y$ are the elements $g\in G$ satisfying $g\cdot x=y$.
Thus $\pi_0(G\ltimes X)=X/G$ (Quillen \cite[\S1]{Quillen73}; in Lean,
\texttt{CategoryTheory.ActionCategory}).

The octonionic world is
\[
  \mathcal H_1=\Gtwo\ltimes\OO^{*}.
\]
Its components are the $\Gtwo$-orbits: the compactified real fixed locus $\RR\cup\{N\}$ and,
for each $\sigma\in\RR$ and $\gamma>0$, the residue-$\CC$ sphere
$S_{(\sigma,\gamma)}=\{\sigma+\gamma v:v\in S^6\}$, labelled by its real centre $\sigma$ and
radius $\gamma$ (Baez \cite{Baez02}; Theorem~\ref{thm:G2-S6}).

The slice-value world $\mathcal S_2$ organizes the same compactified values by their slice
decomposition. Every slice contains the real axis, and
\[
  \OO^{*}=\bigcup_{I\in S^6}\CC_I^{*},\qquad
  \CC_I^{*}\cap\CC_J^{*}=\RR\cup\{N\}\quad(J\neq\pm I).
\]
Its objects are the points of $\OO^{*}$. Its generating paths are the direction arrows induced
by $\Gtwo$ and the band phases $z\mapsto e^{I\theta}z$ within a compactified slice; the Lean
object \texttt{S2} is the quotient path category by the direction-composition relations. In
particular $0$ and $N$ are shared values, not new objects indexed by a zero.

These two geometric views hold the octonionic and slice-value geometry drawn on below.
\end{definition}

\subsection*{The base of the degenerate set and the Grothendieck construction}

\begin{definition}[The projective base, genuine section functor, and total category]\label{def:base}
\lean{GreatCircle.Base, GreatCircle.cayleyGL, GreatCircle.diagonalGL,
GreatCircle.diagonalGL_mul, GreatCircle.diagonalGLHom, GreatCircle.expUnit,
GreatCircle.expUnit_add, GreatCircle.diskExpAction, GreatCircle.diskExpAction_add,
GreatCircle.cayleyProjective, GreatCircle.cayleyProjective_mul,
GreatCircle.orbitRep, GreatCircle.orbitRep_spec, GreatCircle.stabilizerPart,
GreatCircle.orbit_stabilizer_factor, GreatCircle.stabilizerPart_unique,
GreatCircle.stabilizerPart_id, GreatCircle.stabilizerPart_comp,
ASection.distinguishedDiskAction, ASection.distinguishedDiskAction_eq_fullMultiplier,
ASection.distinguishedPoleFactor_euler, ASection.distinguishedPoleFactor_weierstrass,
ASection.eulerPrimeSum, ASection.eulerDiskAction_eq_value,
ASection.GpvTransport, ASection.GpvTransport.diskExpAction_eq_value,
ASection.GpvTransport.diskExpAction_endpoint_eq,
ASection.GpvTransport.lift_endpoint_re_eq, ASection.GpvTransport.level,
ASection.GpvTransport.continuous_level, ASection.GpvTransport.lift_unique,
ASection.GpvTransport.level_independent, ASection.GpvTransport.value_at_source,
ASection.GpvTransport.value_at_target,
ASection.GpvTransport.endpoint_log_norm_eq,
ASection.GpvTransport.endpoint_norm_eq,
ASection.projectiveGpvActionSquare, ASection.projectiveGpvActionSquare_level,
ASection.projectiveGpvActionSquare_input_commutes,
ASection.projectiveGpvActionSquare_output_commutes,
ASection.canonicalAsectionPresentation_euler_prime_stack,
ASection.canonicalAsectionPresentation_euler_lift_exp,
ASection.canonicalAsectionPresentation_euler_lift_unique,
ASection.canonicalAsectionPresentation_euler_level_continuous,
ASection.canonicalAsectionPresentation_euler_winding,
ASection.canonicalAsectionPresentation_euler_toNorth,
ASection.projectiveObjectFrame, ASection.projectiveArrowElement,
ASection.projectiveObjectFrame_north,
ASection.distinguishedDiskAction_fixes_cayley_zero,
ASection.distinguishedDiskAction_fixes_cayley_N,
ASection.projectiveArrowElement_id, ASection.projectiveArrowElement_comp,
ASection.AsectionSlice, ASection.AsectionEquivariant,
ASection.AsectionStateInput, ASection.AsectionStateOutput,
ASection.AsectionState_input_then_equivariant,
ASection.ActionTransportSquare.coordinateTransport_commutes,
ASection.ActionTransportSquare.input_commutes,
ASection.ActionTransportSquare.output_commutes,
ASection.AsectionActionState, ASection.AsectionActionTransport,
ASection.AsectionActionState.ofInput,
ASection.AsectionActionState_ofInput_value,
ASection.AsectionActionTransport_obj_input,
ASection.AsectionActionTransport_obj_positioned,
ASection.AsectionActionTransport_obj_value,
ASection.AsectionActionTransport_id, ASection.AsectionActionTransport_comp,
ASection.AsectionActionDiagram, ASection.residueActionState,
ASection.residueActionState_mem, ASection.c4_infinite,
ASection.transportLevel,
ASection.ambientTotalCategory, ASection.ambientTotalGroupoid,
grothendieckGrpdGroupoid,
CategoryTheory.Grothendieck}
\uses{def:A-section, def:two-worlds, thm:winding-lift, prop:winding-signature}
The base is the projective action groupoid
\[
  \mathcal B
  =PGL(2,\RR)\ltimes\operatorname{OnePoint}(\RR).
\]
Its objects are points of the compactified real projective line and its arrows are projective
transformations together with their transport equations.

The matrix construction makes explicit why the analytic object is simultaneously a function and a
group element. Let
\[
 C=\begin{pmatrix}1&-i\\[2pt]1&i\end{pmatrix},\qquad
 M(z)=\begin{pmatrix}e^z&0\\[2pt]0&1\end{pmatrix}.
\]
The matrix $C$ is the Cayley map $w\mapsto(w-i)/(w+i)$, and $M$ is
exponentiation in matrix form: it takes a complex number $z$ to the diagonal
matrix whose entry is $e^z$. Conjugating an invertible matrix by $C$ and
passing to the projective class is a single further operation, used at every
level of what follows; for $h\in\mathrm{GL}(2,\CC)$ write
\[
  \operatorname{cayleyProjective}(h)=\bigl[C\,h\,C^{-1}\bigr].
\]
The passage from a complex number to a slice M\"obius map is therefore two
steps, exponentiate then conjugate:
\[
  z\;\longmapsto\;M(z)\;\longmapsto\;
  D(z)=\operatorname{cayleyProjective}\bigl(M(z)\bigr)
      \in\operatorname{M\ddot ob}(\CC^*).
\]
Conjugation preserves products, so $\operatorname{cayleyProjective}$ is a
homomorphism; every later appearance of it --- positioning the distinguished
element over the base, and acting by a residual stabilizer factor --- is this
same map applied to a different matrix. Matrix multiplication and
$e^{z+w}=e^ze^w$ give
\[
  M(z+w)=M(z)M(w),\qquad D(z+w)=D(z)D(w).
\]
Thus exponentiation is the homomorphism which turns addition of logarithmic lifts into composition
of the actual slice M\"obius maps. In the Lean source these are the chain
\texttt{diagonalGL}, \texttt{diagonalGLHom}, \texttt{expUnit}, and
\texttt{diskExpAction}, with their multiplication laws. Conditions \textup{C1--C3} select the
A-specific logarithm $\lambda_A$ of the continued nonzero pole factor and hence the one element
\[
  D_A=D(\lambda_A)=\operatorname{cayleyProjective}\bigl(M(\lambda_A)\bigr),
  \qquad
  M(\lambda_A)=\operatorname{diag}\bigl(e^{\lambda_A},1\bigr)
              =\operatorname{diag}(u_A,1).
\]
Here $\lambda_A$ is the logarithm and $M$ is what undoes it: the entry of
$M(\lambda_A)$ is $e^{\lambda_A}=u_A$, the multiplier itself, so $D_A$ is an
exponential action and not a logarithmic one. The logarithm is the coordinate
in which the tape moves; the element it generates lives one exponential above
it.

Which logarithm $\lambda_A$ is, and that it is one at all, is supplied by the
slice facts of Part~2 rather than chosen here. The slice exponential is
slice-preserving and nowhere zero, with $\lvert\exp q\rvert=e^{\operatorname{Re}q}$
\textup{(}Theorem~\ref{thm:slice-exp}\textup{)}; its logarithm manifold
unwraps the multivalued logarithm into a single-valued one, with the branches
carried by the manifold and not by the value
\textup{(}Theorem~\ref{thm:log-manifold}\textup{)}; and a tame lift along a path
is unique once its initial rung is fixed, the ambiguity between branches being
exactly a winding rung in $2\pi i\ZZ$
\textup{(}Theorem~\ref{thm:winding-lift},
Proposition~\ref{prop:winding-signature}\textup{)}.

That ambiguity dies under $M$. Since $e^{\lambda_A+2\pi ik}=e^{\lambda_A}$,
every branch gives the same $M(\lambda_A)$ and hence the same $D_A$: the
logarithm carries the winding, and the exponential action it generates does not
see it. This is what lets the tape run through the whole base while the
distinguished element stays one element, and it is recorded again below as
\textup{(B6)--(B8)}.
The Euler and Weierstrass expressions are the two presentations of this same multiplier. The
function/group-element passage is the commuting dictionary
\[
\begin{tikzcd}[column sep=large,row sep=large]
z \arrow[r,"\exp"] \arrow[d,mapsto] & e^z\in\CC^\times \arrow[d,mapsto]\\
M(z)=\operatorname{diag}(e^z,1)
  \arrow[r,"\operatorname{cayleyProjective}"'] &
D(z)\in\operatorname{M\ddot ob}(\CC^*).
\end{tikzcd}
\]

The common Euler--Weierstrass multiplier can be followed directly through this dictionary. Write
$p_A$ for the pole of $A$ and $g_A=\texttt{distinguishedPoleFactor\,A}$. On the punctured pole chart, conditions
\textup{C2} and \textup{C3} give the two identities
\[
\begin{aligned}
 g_A(z)
   &=(z-p_A)
     \exp\!\left(\sum_{p\in A.\iota}\ell_p(z)\right),\\[2pt]
 g_A(z)
   &=z^{m_A}R_{\mathrm{fac},A}(z)e^{h_A(z)}
     \prod_{n\geq0}\mathcal E\bigl(z;q_{A,n}\bigr).
\end{aligned}
\]
Condition \textup{C1} continues this same factor through the simple pole and proves
$u_A=g_A(p_A)\neq0$. Hence $u_A\in\CC^\times$, its logarithm
$\lambda_A=\log u_A$ satisfies $e^{\lambda_A}=u_A$, and the distinguished element is
\[
 D_A=D(\lambda_A)
   =\operatorname{cayleyProjective}\bigl(\operatorname{diag}(u_A,1)\bigr).
\]
This calculation also gives its two fixed boundary points. In the standard
projective chart $\operatorname{diag}(u_A,1)$ acts as $w\mapsto u_Aw$, which
fixes $0$ and $\infty$ and moves everything else. Conjugation carries a fixed
point to its Cayley coordinate, so the fixed points of $D_A$ are those of
$\operatorname{diag}(u_A,1)$ read through $C$: the projective zero and the
shared north point $N$. Therefore
\[
  D_A(0)=0,\qquad D_A(N)=N,
\]
which are the declarations
\[
\begin{gathered}
\texttt{ASection.distinguishedDiskAction\_fixes\_cayley\_zero},\\
\texttt{ASection.distinguishedDiskAction\_fixes\_cayley\_N}.
\end{gathered}
\]
The Euler expression at $0$ and the
Weierstrass expression at $N$ are consequently the two boundary readings of this one matrix-built
function.

The GPV lift records how that function varies between its boundary presentations. Let
$X,Y\in\mathcal B$. A projective-base transport from $X$ to $Y$ carries three continuous tapes:
a compactified domain path $\delta^*$, a nonzero value path $v$, its compactification $v^*$,
and a logarithmic lift $\Gamma$.
Their defining equations are
\[
\begin{aligned}
 \delta^*(0)&=\operatorname{cayleyCoord}(X),
 &\delta^*(1)&=\operatorname{cayleyCoord}(Y),\\
 v^*(t)&=A^*\bigl(\delta^*(t)\bigr),
 &e^{\Gamma(t)}&=v(t),\\
 \Gamma(1)-\Gamma(0)&=2\pi i k,
 &&k\in\ZZ.
\end{aligned}
\]
These are precisely the fields of \texttt{ASection.GpvTransport}. At each $t$ the lift equation
passes through the matrix dictionary:
\[
\begin{tikzcd}[column sep=huge,row sep=large]
\Gamma(t) \arrow[r,"\exp"] \arrow[d,"D"'] &
v(t)\in\CC^\times
  \arrow[d,"{u\mapsto\operatorname{cayleyProjective}(\operatorname{diag}(u,1))}"]\\
D(\Gamma(t)) \arrow[r,equal] &
\operatorname{cayleyProjective}\bigl(\operatorname{diag}(v(t),1)\bigr).
\end{tikzcd}
\]
Thus every point of the lift is already an action of the same kind as $D_A$. Taking real parts
and norms in the lift equation gives
\[
  \operatorname{Re}\Gamma(t)=\log\lVert v(t)\rVert.
\]
The right-hand side shows that the real level is continuous and depends only on the value tape.
Any second continuous lift with the same initial rung equals $\Gamma$, while every lift of the
same value tape has this same real part. Taking real parts in the winding equation gives
\[
  \operatorname{Re}\Gamma(0)=\operatorname{Re}\Gamma(1),
\]
and exponentiating the full equation gives
\[
  D(\Gamma(0))=D(\Gamma(1)).
\]
The winding integer therefore records the vertical displacement of the logarithmic rung, and the
real level and the endpoint M\"obius action remain fixed.

Condition \textup{C2} now supplies the canonical arithmetic instance of this construction. Along
an Euler half-space loop $\delta$,
\[
  \Gamma_A(t)=\sum_{p\in A.\iota}\ell_p(\delta(t)),
  \qquad e^{\Gamma_A(t)}=A(\delta(t)),
\]
with the prime index retained inside the complete summation. Local normal convergence makes
$\Gamma_A$ continuous, the zero-free half-space makes its value tape lie in $\CC^\times$, and the
GPV uniqueness theorem fixes the whole lift once $\Gamma_A(0)$ is chosen. Since $\delta(0)=
\delta(1)$, the same prime sum occurs at both endpoints, so its checked winding is $k=0$. At every
instant the square \texttt{canonicalAsectionPresentation\_euler\_toNorth} carries this complete
prime tape to the common north frame. Condition \textup{C3} supplies the Weierstrass boundary
presentation there, and condition \textup{C1} identifies both presentations with the continued
factor $g_A$. The prime sum, its value under $A$, its logarithmic level and winding, and the matrix
action it generates are therefore simultaneous coordinates of the distinguished element.

The projective action is read in the direction-indexed sphere-world groupoid. Its objects are
directions $I\in S^6$ equipped with their compactified slice spheres, and its morphisms carry the
direction and M\"obius legs. A normalized value state pairs such a direction with a point of its
slice.

We now have to put $D_A$ at every object of the base, and here the choice of an
action groupoid pays for itself: no chart is required. In an action groupoid an
arrow \emph{is} a group element, and two objects are joined exactly when one
carries the other. So the base is a single orbit, every object is reached from
the north object $N$ by some element, and whatever $D_A$ becomes at another
object is its image under that element --- an element of the orbit of $D_A$,
not a new object needing coordinates to describe. The only freedom is which
element does the carrying, and fixing that once removes it.

Fix, then, the canonical orbit representative $o_b:N\to b$, one element carrying
$N$ to $b$, and define the frame at $b$ to be $D_A$ carried along it:
\[
  a_A(b)=\operatorname{cayleyProjective}(o_b)\,D_A .
\]
At the north object itself the representative is the identity, $o_N=1$, so
\[
  a_A(N)=\operatorname{cayleyProjective}(1)\,D_A=D_A .
\]
The frame at north \emph{is} the distinguished element, uncarried --- which is
why every reading taken at $N$ later in this paper is a reading of $D_A$ itself,
and why the residue condition can be stated there without a chart.

With the representatives fixed, an arbitrary arrow $f:x\to y$ factors through
$N$ in exactly one way: go back along $o_x$, do something that fixes $N$, and
come out along $o_y$. That middle part is the residual north-stabilizer factor
\[
  r_f=o_y^{-1}f\,o_x,\qquad
  f=o_y r_f o_x^{-1},
\]
and it is the only one, since $o_x,o_y$ are fixed and invertible: $f=o_yro_x^{-1}$
determines $r=o_y^{-1}fo_x$ outright. Uniqueness of the stabilizer part is
therefore the invertibility of the representatives, and every later appeal to it
--- including at the two boundary transports of the comparison in
Lemma~\ref{lem:c-residue-transitive} --- is this one line
\textup{(}\texttt{GreatCircle.stabilizerPart\_unique}\textup{)}.

Call this assignment $a_A$: it gives each object $b$ the frame $a_A(b)$ above,
and each arrow $f:x\to y$ the M\"obius element $a_A(f)$ carrying the frame at
$x$ to the frame at $y$, namely its residual factor conjugated between the two
frames. Together,
\[
  a_A(b)=\operatorname{cayleyProjective}(o_b)\,D_A,\qquad
  a_A(f)=a_A(y)\,\operatorname{cayleyProjective}(r_f)\,a_A(x)^{-1},
                              \tag{OS}\label{eq:orbit-stabilizer}
\]
where $f\fatsemi g$ denotes composition in diagrammatic order: $f$ first, then
$g$. Since the representatives are fixed, $r_{1_x}=1$ and
$r_{f\fatsemi g}=r_g r_f$; since $\operatorname{cayleyProjective}$ is a
homomorphism, those two identities are carried to $a_A$, which is what makes
the identity and composition laws hold.

Equation \eqref{eq:orbit-stabilizer} is the orbit--stabilizer well-definedness
of every functor in this paper, and it is the only device used to obtain one.
Each is built the same way: the one element $D_A$, carried to every object along
the fixed representatives, and to every arrow by conjugating its residual factor
between the frames at the arrow's two ends.
We refer to \eqref{eq:orbit-stabilizer} by name throughout, and every later
positioning statement is an instance of it rather than a new hypothesis.

The preceding construction gives nine checkable readings on the projective base:
\begin{enumerate}[label=\textup{(B\arabic*)},leftmargin=3.2em,itemsep=2pt]
\item the C2 Euler and C3 Weierstrass formulas present the one C1-continued factor $g_A$;
\item $u_A=g_A(p_A)$ and $\lambda_A=\log u_A$ generate the one matrix action $D_A$;
\item the diagonal calculation fixes the two boundary points $0$ and $N$;
\item the Euler lift retains the complete prime-indexed sum
  $\sum_{p\in A.\iota}\ell_p$;
\item exponentiation identifies the lift pointwise with the A-value and with its M\"obius matrix;
\item the initial rung determines the unique tame continuous lift;
\item its real level is $\log\lVert A\rVert$, continuous and independent of the lift branch;
\item its winding lies in $2\pi i\ZZ$, preserves the endpoint action and real level, and equals
  zero on the C2 Euler loop;
\item the unique orbit and stabilizer factors carry every instant of this action through the base,
  with identity and composition inherited from their factorization laws.
\end{enumerate}
Each line is a consequence of the construction above: \textup{(B1)--(B3)} come from the common
pole factor and its matrix, \textup{(B4)--(B8)} from its GPV lift, and \textup{(B9)} from the
projective orbit--stabilizer sweep.

The direction-and-M\"obius projection is
\[
  A^{\mathrm{slice}}:\mathcal B\longrightarrow\SphereWorld,
\]
written \texttt{AsectionSlice\,A} in the Lean source. It records the slice sphere and its
M\"obius transport. The full A-section functor below additionally retains the normalized point
and the value produced on that sphere.

The slice-preserving realization on $\Ostar$ is the $G_2$-equivariant endofunctor
$A_{\OO}:\mathcal H_1\to\mathcal H_1$.

The states it acts on form the \emph{normalized state world}
$\operatorname{StateWorld}_A$ \textup{(}\texttt{AsectionStateWorld\,A}\textup{)}:
the action groupoid of $G_2$ on normalized A-section states. An object is a
slice direction $I\in S^6$ together with a normalized coordinate $u$ in the
slice that direction selects --- a point of the sphere $A^{\mathrm{slice}}$
records, now carrying that point and not only the sphere. An arrow is a single
element $\gamma\in G_2$, moving the direction and carrying the coordinate with
it. The state world has two functors to $\mathcal H_1$, reading a state's
coordinate before and after $A_{\OO}$ acts: the input eye
$\operatorname{In}_A$ and the output eye $\operatorname{Out}_A$.
The round trip is the checked triangle
\[
\begin{tikzcd}[column sep=large,row sep=large]
\operatorname{StateWorld}_A
  \arrow[r,"\operatorname{In}_A"]
  \arrow[dr,"\operatorname{Out}_A"'] &
\mathcal H_1 \arrow[d,"A_{\OO}"]\\
& \mathcal H_1,
\end{tikzcd}
\qquad
\operatorname{Out}_A=\operatorname{In}_A\fatsemi A_{\OO}.
\]
This equality is
\texttt{AsectionState\_input\_then\_equivariant}: it holds on arrows as well as objects, so
the normalized input coordinates genuinely drive the realized A-values.

Positioning this round trip over $b\in\mathcal B$ forms a constrained graph. An object in the
fibre consists of an input state $x$, the forced positioned state
$a_A(b)\cdot x$, and the forced value
$\operatorname{Out}_A(a_A(b)\cdot x)$. Thus both the positioned point and its value are determined
by the input and the action. For $f:b\to b'$, the two legs of the orbit--stabilizer square transport the input
and positioned faces, and its already-proved commuting equation makes the target graph
well-defined. Reading that same square through $\operatorname{In}_A$ and
$\operatorname{Out}_A$ gives the checked input and output naturality squares. Identity and
composition of these graph transports are inherited from the unique factorization of $r_f$.
Sweeping the whole constrained graph through the base therefore gives the genuine A-section
functor
\[
  \mathcal A_A:\mathcal B\longrightarrow\mathbf{Grpd},
\]
written \texttt{AsectionActionDiagram\,A} in the Lean source.

Both its source and its targets are action groupoids. The source
$\mathcal B=\mathrm{PGL}(2,\RR)\ltimes\mathrm{OnePoint}(\RR)$ is the projective
action groupoid: an object is a point of the compactified real projective line,
and an arrow $f:x\to y$ is a projective transformation together with its
transport equation, so it is one element of $\mathrm{PGL}(2,\RR)$ carrying $x$
to $y$. Each groupoid it produces is an action groupoid of $G_2$ on states, and
the element positioned there is the distinguished one, $D_A$.

On a base object $b$, the groupoid $\mathcal A_A(b)$ consists of the normalized
A-section states carrying $D_A$ to that object. The carrying is done by the
object frame $a_A(b)=\operatorname{cayleyProjective}(o_b)D_A$, which is $D_A$
conjugated along the canonical orbit representative $o_b$ --- the same element
at every base object, moved into position. An object of
$\mathcal A_A(b)$ is one constrained graph \eqref{eq:north-graph}: a slice
direction $I\in S^6$ and a normalized input coordinate $u$, together with the
positioned coordinate $a_A(b)\mathbin{\cdot}(I,u)$ that this frame determines
and the A-value realized there. An arrow is a single element $\gamma\in G_2$
acting on the slice direction and carrying the rest of the graph with it, so
$\mathcal A_A(b)$ is the action groupoid of $G_2$ on those states.

On a base arrow $f:x\to y$ --- that is, on one projective transformation ---
the value $\mathcal A_A(f)$ is transport by the arrow element
\[
  a_A(f)=a_A(y)\,\operatorname{cayleyProjective}(r_f)\,a_A(x)^{-1}
\]
of \eqref{eq:orbit-stabilizer}: the distinguished action read at the target
frame, composed with the Cayley action of the residual factor $r_f$ that $f$
leaves after its orbit part is stripped, composed with the inverse of the
distinguished action read at the source frame. On a state this moves the input
coordinate by the input leg $\operatorname{cayleyProjective}(r_f)$ and the
positioned coordinate by the positioned leg, the two horizontal legs of the
square whose verticals are those two frames. On an arrow it is the identity on
$\gamma$, since $G_2$ acts on the slice direction while $a_A(f)$ acts on the
coordinate, so the two commute. The identity and composition laws hold because
the fixed representatives force $r_{1_x}=1$ and
$r_{f\fatsemi g}=r_g r_f$, and $\operatorname{cayleyProjective}(-)$ is a
homomorphism.

The nine base readings \textup{(B1)--(B9)} pass through this constrained graph by the certified
projective/GPV action squares. Three further readings are thereby realized in the generated
states. First, condition \textup{C3} supplies the populated residue-$\CC$ sphere states, represented
by \texttt{residueActionState} and certified by \texttt{residueActionState\_mem}. Second,
condition \textup{C4} makes their population infinite through \texttt{c4\_infinite}. Third, every
such state carries its intrinsic real coordinate through \texttt{transportLevel}. These are the
three output readings. Together with \textup{(B1)--(B9)} they form the $9+3$ partition: nine
properties of the one distinguished action on the base, followed functorially by three readings
of the A-generated residue states in the total.

The total value-transport category is
\[
  \mathcal T_A=\int_{\mathcal B}\mathcal A_A.
\]
An object is a base position together with a normalized A-section value state in its fibre. A
morphism is a base channel together with the compatible A-section transport supplied by $A$.
The populated residue-$\CC$ zero-sphere states are the degenerate-fibre outputs of this
completed construction, produced by \textup{C3--C4}. The associated diagram
$\pi_0\circ\mathcal A_A:\mathcal B\to\mathbf{Set}$ is the components diagram required by the
Grothendieck/colimit readout (Lemma~\ref{lem:pi0-grothendieck}).
\end{definition}


\section{The concentricity theorem}\label{sec:detector}

We now have the worlds and the base, so this section defines the class of functions precisely and
proves the theorem. Two facts about slice-preserving functions come first: the Weierstrass
factorization that gives condition C3 its meaning, and the fact that a residue-$\CC$ zero is a whole
$6$-sphere, a single connected component of the octonionic world. Then the categorical engine,
Riehl's computation of the components of a Grothendieck construction as a colimit, which sorts the
assembled total into the classes linked by the section's real-value transports. With these in hand
we define an A-section by its four conditions C1--C4, prove that the C-residue total of any
A-section is transitive with a singleton component colimit, identify exactly what that singleton
does and does not give, and exhibit an A-section whose zero-spheres are not concentric.

\begin{proposition}[Slice-regular Weierstrass factorization; the content of C3]\label{prop:weierstrass}
\lean{weierstrassE, spherePrimary, ASection.c3_multipliable,
ASection.c3_locMajorant, ASection.c3_factorization}
\uses{def:R}
Let $A\in\mathcal R$ have a zero of order $m\ge 0$ at the origin and no non-real isolated zeros
\textup{(}for $A\in\mathcal R$ the latter is automatic: non-real zeros come in whole spheres,
Lemma~\ref{lem:residue-spheres}\textup{)}. Then on $\OO^{*}\setminus\{N\}$ it factors as
\[
  A=q^{m}\,R\,S\,h,
\]
where $q$ denotes the octonionic coordinate function \textup{(}the identity section, so $q^m$
carries the order-$m$ zero at the origin\textup{)}, $R$ is a slice-preserving factor vanishing
exactly at the real zeros of $A$, $S$ is a factor vanishing exactly at its spherical
\textup{(}residue-$\CC$\textup{)} zero-spheres, and $h$ is never vanishing; each of $R,S$ is the
slice-regular Weierstrass product over the corresponding zeros. The factorization holds
\emph{including the case where either family of zeros is infinite}.
\end{proposition}
\begin{proof}
This is the factorization theorem of Altavilla--de Fabritiis \cite[Prop.~3.1, Thm.~3.2]{AdFslice},
whose primary factors are the slice-regular ones of Gentili--Vignozzi \cite{GV}; convergence of
the infinite products is part of the cited statement. The pair $R,S$ is the residue split named
in \textup{C3}: $R$ carries the residue-$\RR$ \textup{(}classically trivial\textup{)} zeros and
$S$ the residue-$\CC$ \textup{(}classically nontrivial\textup{)} ones, the latter infinite in
number by \textup{C4}.
\end{proof}

\begin{lemma}[Residue-$\CC$ zeros are $6$-sphere components of $\mathcal H_1$]\label{lem:residue-spheres}
\lean{ASection.sphereZero_mem_CResidueZeroLocus,
ASection.mem_CResidueZeroLocus_iff_exists_sphereZero,
G2.exists_smul_eq_of_mem_unitImaginarySphere}
\uses{def:two-worlds, thm:G2-S6, thm:wang}
For $A\in\mathcal R$, the non-real zeros of $A$ are $6$-spheres, each a single connected
component of the translation groupoid $\mathcal H_1$.
\end{lemma}
\begin{proof}
$A$ is slice-preserving, so on a sphere $\sigma+\gamma S^6$ \textup{(}$\gamma>0$\textup{)} it has
the form $A(\sigma+\gamma v)=\alpha(\sigma,\gamma)+v\,\beta(\sigma,\gamma)$ with $\alpha,\beta$
real (the real stem, \cite{Wang}); since $v$ is an imaginary unit, $\alpha+v\beta=0$ forces
$\alpha=\beta=0$, independently of $v$. Hence a non-real zero occurs on a \emph{whole} sphere
$\sigma+\gamma S^6$ or at none of its points (the spherical-zero structure, \cite{AdF}); real
zeros are isolated real points (residue field $\RR$). Each such sphere $\sigma+\gamma S^6$ is the
$\Gtwo$-orbit of any of its points (Theorem~\ref{thm:G2-S6}: $\Gtwo$ acts transitively on $S^6$
with stabiliser $\mathrm{SU}(3)$), and the connected components of $\mathcal H_1$ are exactly the
$\Gtwo$-orbits: in a groupoid two objects share a component iff a morphism joins them, and a
morphism $x\to y$ of $\mathcal H_1$ is a $g\in\Gtwo$ with $g\cdot x=y$. Therefore the
residue-$\CC$ zeros, the non-real ones, are exactly the $6$-sphere components of
$A^{-1}(0)$.
\end{proof}

\begin{lemma}[$\pi_0$ of a Grothendieck construction]\label{lem:pi0-grothendieck}
\lean{CategoryTheory.Grothendieck, pi0GrothendieckEquiv, pi0_grothendieck}
\uses{def:base}
For a functor $F:\mathcal B\to\Grpd$, the connected-components functor carries the Grothendieck
construction to the colimit of the component diagram:
\[
  \pi_0\Bigl(\int_{\mathcal B}F\Bigr)\ \cong\ \operatorname*{colim}_{\mathcal B}\,(\pi_0\circ F).
\]
\end{lemma}
\begin{proof}
By Riehl's proof of Lemma~8.3.4, for any $X:\mathcal C\to\mathbf{Set}$ each arrow of the
category of elements imposes exactly the condition that the corresponding elements be identified
in every cone, so that $\pi_0(\operatorname{el}X)\cong\operatorname{colim}_{\mathcal C}X$
\cite[p.~102]{Riehl}. Taking $X=\pi_0\circ F$, the colimit performs exactly the identifications
generated by the maps of $F$. In the Lean development this equivalence is
\texttt{pi0GrothendieckEquiv}, built from the canonical cocone \texttt{pi0Cocone} and
\texttt{colimit.desc} over Mathlib's \texttt{colimit\_sound} and \texttt{colimit\_eq}.
\end{proof}

\begin{definition}[$A$-sections]\label{def:A-section}
\lean{ASection, weierstrassE, spherePrimary}
\uses{def:R, prop:weierstrass, lem:residue-spheres}
An \emph{$A$-section} is a section $A$ of the ring $\mathcal R$ of slice-preserving
slice-regular functions on $\OO^{*}=S^8$ \textup{(Definition~\ref{def:R})} with the following
four properties.
\begin{enumerate}[label=\textup{(C\arabic*)},leftmargin=3.2em]
\item \emph{Meromorphic continuation; single simple pole.} $A$ is meromorphically continued to
  $\OO^{*}\setminus\{N\}$ with exactly one pole; it is simple, lies at a real point, and its
  value there is the point at infinity $\infty=N$. At the domain point $N$ the section carries
  an assigned value in the target sphere, a real point or $\infty$, as data; no continuity or
  meromorphy at $N$ is required, and none holds for $\zetaO$
  \textup{(}Definition~\ref{def:zeta-Cstar}\textup{)}.
\item \emph{Infinite Euler product.} On a slice right half-space $\Omega_0\subset\OO^{*}$,
  $A=\exp\!\big(\sum_p\ell_p\big)$ for an \emph{infinite} family $\{\ell_p\}$ of
  slice-preserving slice-regular functions each zero-free on $\Omega_0$, the family summable
  \emph{locally normally} on $\Omega_0$ \textup{(}the sense in which the cited Euler products
  converge; Theorem~\ref{thm:euler}\textup{)}; in particular $A$ is
  zero-free on $\Omega_0$.
\item \emph{Infinite Weierstrass factorization.} Over the \emph{full divisor} of $A$,
  including the single pole of \textup{(C1)}, at the real point $p_0$, multiplicity $-1$,
  carried by the explicit left factor, on $\OO^{*}\setminus\{p_0\}$,
  \[
    (q-p_0)\,A\;=\;q^{m}\,R\;e^{g}\prod_n\mathcal E(\,\cdot\,;q_n),
  \]
  where $q$ is the octonionic coordinate, $m\ge0$, $R$ is the slice-regular Weierstrass product
  over the \emph{nonzero} residue-$\RR$ \textup{(}real\textup{)} zeros of $A$, vanishing only
  on $\RR$, $g$ is slice-preserving entire, $\{q_n\}$ enumerates the residue-$\CC$ zero-spheres
  of $A$, and $\mathcal E(\,\cdot\,;q_n)$ is the slice-regular Weierstrass primary factor of
  $q_n$ \textup{(}Proposition~\ref{prop:weierstrass}; \cite[Prop.~3.1, Thm.~3.2]{AdFslice},
  \cite{GV}\textup{)}; the factorization holds with either family infinite, the product
  converging \emph{locally normally} on $\OO^{*}\setminus\{p_0\}$ \textup{(}the convergence of
  Proposition~\ref{prop:weierstrass}\textup{)}. \textup{(}The left
  factor makes ``full divisor'' literal: the divisor of a meromorphic section includes its
  single pole; classically, Hadamard factors $(s-1)\zeta(s)$.\textup{)}
\item \emph{Infinitely many residue-$\CC$ zeros.} The index set $\{q_n\}$ is infinite. A closed
  zero of such a section is a real point \textup{(}residue field $\RR$\textup{)} or a $6$-sphere
  $S_{(\sigma,\gamma)}$ \textup{(}residue field $\CC$; \cite{AdFslice}\textup{)}; residue-$\CC$
  names the classically nontrivial zeros \textup{(}Part~2\textup{)}.
\end{enumerate}
\end{definition}

\begin{definition}[The residue subdiagram and its inclusion]\label{def:residue-subdiagram}
\lean{ASection.CResidueZeroLocus, ASection.IsCResidueInput,
ASection.CResidueInputWorld, ASection.AsectionCResidueInputDiagram,
ASection.AsectionCResidueDiagram, ASection.AsectionCResidueInclusion,
ASection.AsectionCResidueInclusionTotal,
ASection.AsectionCResidueInclusionTotal_full,
ASection.AsectionCResidueInclusionTotal_faithful,
ASection.residueTotalCategory, ASection.residueTotalGroupoid}
\uses{def:base, def:A-section, lem:residue-spheres}
The \emph{residue locus} of an A-section is $\{z : A(z)=0,\ \operatorname{Im}z>0\}$
\textup{(}\texttt{CResidueZeroLocus}\textup{)}. The \emph{residue subdiagram}
$\mathcal R_A:\mathcal B\to\mathbf{Grpd}$ \textup{(}\texttt{AsectionCResidueDiagram}\textup{)} is
the preimage of the states this locus selects at the north frame --- a preimage taken in the
total action groupoid, not of a set --- with the witnessing base arrow carried as data.

It has the same source as $\mathcal A_A$, the projective action groupoid
$\mathcal B$, and its values are again action groupoids of $G_2$ --- the
selection cuts down which states are present without altering what acts on
them.

On a base object $X$, the value $\mathcal R_A(X)$ is the full subgroupoid of
$\mathcal A_A(X)$ on the witnessed states. An object there is a state $x$ of
$\mathcal A_A(X)$ together with a north residue state $x_N$ and a base arrow
$g:N\to X$ --- one projective transformation carrying the north point to $X$
--- satisfying $\mathcal A_A(g)(x_N)=x$; the witness travels with the object.
An arrow is an arrow of $\mathcal A_A(X)$ between two such states, so again a
single element $\gamma\in G_2$: the selection is on objects, and $G_2$ still
acts.

On a base arrow $f:X\to Y$, the value $\mathcal R_A(f)$ is the restriction of
the ambient transport $\mathcal A_A(f)$, which is transport by the arrow
element $a_A(f)$ carrying $D_A$ between the frames at $X$ and $Y$. That
restriction lands in the subgroupoid because a witnessed state stays witnessed,
with the \emph{same} north state and the composite witness $g\fatsemi f$:
functoriality of $\mathcal A_A$ gives
$\mathcal A_A(g\fatsemi f)(x_N)=\mathcal A_A(f)\bigl(\mathcal A_A(g)(x_N)\bigr)
 =\mathcal A_A(f)(x)$.
So the north state is what is preserved, and the witnessing projective
transformation is what accumulates.
Its inclusion $\iota_A:\mathcal R_A\Rightarrow\mathcal A_A$
\textup{(}\texttt{AsectionCResidueInclusion}\textup{)} sends a witnessed object
to its ambient state and is the identity on arrows; it is fibrewise fully
faithful onto that inverse image, and its naturality squares retain the ambient
A-section transport itself.
\end{definition}

\begin{lemma}[Transitivity of the C-residue total]\label{lem:c-residue-transitive}
\lean{ASection.CResidueInputActionSquare,
ASection.CResidueInputFiberHom,
ASection.CResidueInputTotal_transitive,
ASection.sweepTransitive_on_residueSystem}
\leanok
\uses{def:A-section, def:base, def:residue-subdiagram, prop:weierstrass,
lem:residue-spheres, thm:G2-S6}
For an A-section $A$, the total inverse-image groupoid
\[
  \int_{\mathcal B}\mathcal R_A
\]
is transitive: any two of its objects are joined by a morphism.
\end{lemma}
\begin{proof}
Let
\[
  P,Q\in\int_{\mathcal B}\mathcal R_A
\]
be arbitrary. By the inverse-image definition, the fibre data of $P$ and $Q$
supply north residue states $x_N,y_N$ together with base arrows
\[
  g:N\longrightarrow P.\mathrm{base},\qquad
  h:N\longrightarrow Q.\mathrm{base},
\]
whose A-section transports recover the two given fibre states. So the arbitrary
problem reduces to a north endomorphism $k:N\to N$ and a fibre morphism
\[
  \varphi:\mathcal A_A(k)(x_N)\longrightarrow y_N .
\]
Once these are in hand, functorial transport along
$g^{-1}\fatsemi k\fatsemi h$ gives an ambient Grothendieck morphism $P\to Q$,
and the full inclusion $\int\iota_A$ returns that morphism to
$\int_{\mathcal B}\mathcal R_A$. It remains only to construct the north
comparison.

\medskip
\noindent\emph{The state is a triple.}
At $N$ the A-section action functor records
\[
  (I,u)\longmapsto
  \bigl(\,\underbrace{(I,u)}_{\text{input}},\;
        \underbrace{D_A\mathbin{\cdot}(I,u)}_{\text{positioned}},\;
        \underbrace{A_{\OO}(D_A\mathbin{\cdot}(I,u))}_{\text{realized value}}\,\bigr).
                                    \tag{A}\label{eq:north-graph}
\]
The three entries carry the three things the comparison needs, and they carry
them separately. The residue condition of $\mathcal R_A(N)$ selects on the
\emph{middle} entry: a north residue state is one whose positioned coordinate is
an authored C3 coordinate. That is the entry the inverse-image construction is
built on.

\medskip
\noindent\emph{The two states, and where they differ.}
Write the two supplied states as evaluations of \eqref{eq:north-graph}:
\[
\begin{aligned}
  x_N&=\bigl((I_1,u_1),\;z_1,\;A_{\OO}(z_1)\bigr),
    &z_1&=D_A\mathbin{\cdot}u_1,\\[0.5ex]
  y_N&=\bigl((I_2,u_2),\;z_2,\;A_{\OO}(z_2)\bigr),
    &z_2&=D_A\mathbin{\cdot}u_2,
\end{aligned}
                                    \tag{P}\label{eq:north-preimages}
\]
with $z_1,z_2$ the authored C3 coordinates their residue condition selects and
$I_1,I_2\in S^6$ their slice directions. Since $D_A$ is invertible, each middle
entry has exactly one input behind it, so \eqref{eq:north-preimages} determines
$u_1$ and $u_2$ rather than choosing them. The two states differ in their left
entries and in nothing else.

These inputs sit at $N$. They are not the input the two runs are read at: that
one sits at the source of the tape and is written $u_\ast$, and the two runs
carry it to $u_1$ and to $u_2$ respectively.

\medskip
\noindent\emph{What they share.}
Each state is one sweep of the one construction, and each sweep runs the fixed
$0$-to-$N$ tape once:
\[
  \operatorname{eulerToNorth}_1,\ \operatorname{eulerToNorth}_2
    \;:\;0\longrightarrow N .
\]
Along that tape condition \textup{C2} supplies the arithmetic instance
\[
  \Gamma_A(t)=\sum_{p\in A.\iota}\ell_p(\delta(t)),
  \qquad e^{\Gamma_A(t)}=A(\delta(t)),
\]
with the prime index retained inside the complete
summation. Local normal convergence makes $\Gamma_A$ continuous, the zero-free
half-space keeps its value tape in $\CC^\times$, and GPV uniqueness fixes the
whole lift once $\Gamma_A(0)$ is chosen; since $\delta(0)=\delta(1)$ the same
prime sum occurs at both endpoints, so the checked winding is $k=0$. At every
instant the \texttt{euler\_toNorth} square of the canonical presentation
carries this complete prime tape to the common north frame, where \textup{C3}
supplies the Weierstrass boundary presentation and \textup{C1} identifies both
presentations with the continued factor $g_A$.

Euler presents at $0$, and the same
presentation arriving at $N$ \emph{is} Weierstrass. That identity is what makes
the common datum nonzero, $u_A=g_A(p_A)\neq0$
\textup{(}Definition~\ref{def:base}\textup{)}, hence $D_A$ invertible, hence
\eqref{eq:orbit-stabilizer} available at every instant of the tape up to and
including $N$. Invertibility at $N$ is not an assumption imposed on $D_A$; it is
what the continuation delivers at the point where the two presentations agree.

\medskip
\noindent\emph{The two sweeps as squares.}
Write
\[
  r_1=\operatorname{stabilizerPart}(\operatorname{eulerToNorth}_1),\qquad
  r_2=\operatorname{stabilizerPart}(\operatorname{eulerToNorth}_2)
\]
for their residual factors, and write
\[
  R_i=\operatorname{cayleyProjective}(r_i)
\]
for the Cayley action each one generates. The $r_i$ are forced, not chosen:
orbit--stabilizer gives
\[
  \operatorname{eulerToNorth}_i=o_N\,r_i\,o_0^{-1},
\]
and the fixed representatives determine the stabilizer part of that
factorization outright, as shown at \eqref{eq:orbit-stabilizer}.

Read each sweep as a commuting square. The verticals are the object frames
already fixed --- $a_A(0)$ at the source, $a_A(N)=D_A$ at north. The horizontals
are the two legs of \eqref{eq:orbit-stabilizer}: on top the input leg $R_i$,
along the bottom the arrow element
$a_A(\operatorname{eulerToNorth}_i)$. Read at $x=0$, $y=N$,
\eqref{eq:orbit-stabilizer} says each square commutes:
\[
  a_A(\operatorname{eulerToNorth}_i)\,a_A(0)=D_A\,R_i .
                                    \tag{S}\label{eq:north-face-squares}
\]
Equation \eqref{eq:north-face-squares} is \eqref{eq:orbit-stabilizer} read at
the base component of each sweep, so it is derived from the positioning already
fixed. The two sweeps are therefore
\[
\begin{tikzcd}[column sep=3em,row sep=3.2em,ampersand replacement=\&]
u_\ast \arrow[r,"R_1"] \arrow[d,"a_A(0)"']\&
u_1 \arrow[d,"D_A"]\\
a_A(0)\mathbin{\cdot}u_\ast
  \arrow[r,"a_A(\operatorname{eulerToNorth}_1)"']\&
z_1
\end{tikzcd}
\]
\[
\begin{tikzcd}[column sep=3em,row sep=3.2em,ampersand replacement=\&]
u_\ast \arrow[r,"R_2"] \arrow[d,"a_A(0)"']\&
u_2 \arrow[d,"D_A"]\\
a_A(0)\mathbin{\cdot}u_\ast
  \arrow[r,"a_A(\operatorname{eulerToNorth}_2)"']\&
z_2
\end{tikzcd}
\]
one square apiece, sharing their top-left corner.

\medskip
\noindent\emph{Both sweeps read at the one source input.}
The two runs leave the same source object, so they are read at its normalized
input coordinate $u_\ast$. The C3 boundary readings identify the positioned
outputs with the authored coordinates,
\[
  a_A(\operatorname{eulerToNorth}_i)\mathbin{\cdot}
    \bigl(a_A(0)\mathbin{\cdot}u_\ast\bigr)=z_i ,
                                    \tag{B}\label{eq:north-boundary-readings}
\]
and \eqref{eq:north-preimages} rewrites each as $z_i=D_A\mathbin{\cdot}u_i$.
Evaluating \eqref{eq:north-face-squares} at $u_\ast$, reading the left side by
\eqref{eq:north-boundary-readings} and the right by
\eqref{eq:north-preimages}:
\[
  D_A\mathbin{\cdot}\bigl(R_i\mathbin{\cdot}u_\ast\bigr)
    =a_A(\operatorname{eulerToNorth}_i)\mathbin{\cdot}
      \bigl(a_A(0)\mathbin{\cdot}u_\ast\bigr)
    =z_i=D_A\mathbin{\cdot}u_i .
\]
$D_A$ is invertible, so cancelling it on the left gives the two input equations
\[
  R_1\mathbin{\cdot}u_\ast=u_1,\qquad
  R_2\mathbin{\cdot}u_\ast=u_2 .
                                    \tag{I}\label{eq:north-input-readings}
\]
Both sweeps are now read at the one source input, and their difference is
exactly the difference of their residual factors.

\medskip
\noindent\emph{The coordinate half.}
Put $r=r_2r_1^{-1}$, the residual factor of the second sweep relative to the
first. Then \eqref{eq:north-input-readings} gives
\[
\begin{aligned}
  \operatorname{cayleyProjective}(r)\mathbin{\cdot}u_1
    &=\operatorname{cayleyProjective}(r_2r_1^{-1})\,
       R_1\mathbin{\cdot}u_\ast
     =R_2\mathbin{\cdot}u_\ast=u_2,\\[0.6ex]
  D_A\,\operatorname{cayleyProjective}(r)\,D_A^{-1}\mathbin{\cdot}z_1
    &=D_A\mathbin{\cdot}
       \bigl(\operatorname{cayleyProjective}(r)\mathbin{\cdot}u_1\bigr)
     =D_A\mathbin{\cdot}u_2=z_2 .
\end{aligned}
\]
The corresponding base loop is
\[
  k=\operatorname{eulerToNorth}_1^{-1}
    \fatsemi\operatorname{eulerToNorth}_2\;:\;N\longrightarrow N,
\]
and since composing stabilizer parts multiplies them in the same order as the
action while inversion inverts them,
\[
  \operatorname{stabilizerPart}(k)=r_2r_1^{-1}=r .
                                    \tag{R}\label{eq:north-relative-stabilizer}
\]
Equivalently, from $\operatorname{eulerToNorth}_i=o_N\,r_i\,o_0^{-1}$: the first
sweep inverted and followed by the second cancels the common $0$-frame $o_0$,
and $o_N=1$ at north. So the input leg of $\mathcal A_A(k)$ carries $u_1$ to
$u_2$, and conjugating by $D_A$ carries $z_1$ to $z_2$.

\medskip
\noindent\emph{The direction half, and why it is available here.}
The coordinate half moves the middle entry. The right entry still has to move,
and that is the entry the sweep produced.

What $\mathcal A_A$ sweeps is the slice, and what indexes the slices is the
direction: $I$ ranges over $S^6$, the unit imaginary octonions, $D_A$ positions
the coordinate in the slice that $I$ selects, and $A_{\OO}$ realizes the value
there. So the directions carried by $x_N$ and $y_N$ are not labels attached
afterwards --- they are the indices the functor swept over, and every one of
them is a unit imaginary octonion because that is what a slice direction is
\textup{(}Definition~\ref{def:carrier}\textup{)}. This is exactly the set on
which $G_2$ is transitive \textup{(}Theorem~\ref{thm:G2-S6}\textup{)}, so it
supplies $\gamma\in G_2$ with
\[
  \gamma I_1=I_2 .
\]
The ambient total cannot take this step. Its objects are all the states of
$\mathcal A_A$, and $G_2$ moves a state only within the sphere its direction
already lies on. The inverse image, by contrast, is taken over precisely what
the sweep produces --- the residue states on the C3 zero spheres --- so the
transitivity of $G_2$ on the unit imaginary octonions acts on the whole of what
$\int_{\mathcal B}\mathcal R_A$ contains, and on nothing larger. That is why
this comparison is available there and not in $\int_{\mathcal B}\mathcal A_A$.

The transported coordinate equality and this direction action are the two
components of an arrow
\[
  \varphi:\mathcal A_A(k)(x_N)\longrightarrow y_N .
                                    \tag{$\Phi$}\label{eq:north-fibre-arrow}
\]

\medskip
\noindent\emph{Return to $P$ and $Q$.}
Their inverse-image data supply the arrows $g,h$ above together with
\[
  \mathcal A_A(g)(x_N)=P.\mathrm{fiber},\qquad
  \mathcal A_A(h)(y_N)=Q.\mathrm{fiber},
                                    \tag{N}\label{eq:north-witness-readings}
\]
and the base component of the required total arrow is
\[
  (g^{-1}\fatsemi k)\fatsemi h:
  P.\mathrm{base}\longrightarrow Q.\mathrm{base}.
                                    \tag{G}\label{eq:residue-total-base-arrow}
\]
A morphism $P\to Q$ of $\int_{\mathcal B}\mathcal A_A$ is a base arrow
$f:P.\mathrm{base}\to Q.\mathrm{base}$ together with a fibre arrow
$\mathcal A_A(f)(P.\mathrm{fiber})\to Q.\mathrm{fiber}$. Functoriality gives
$\mathcal A_A\bigl((g^{-1}\fatsemi k)\fatsemi h\bigr)
 =\mathcal A_A(h)\circ\mathcal A_A(k)\circ\mathcal A_A(g^{-1})$, and
$\mathcal A_A(g^{-1})(P.\mathrm{fiber})=x_N$ by
\eqref{eq:north-witness-readings}. So the composite carries
$P.\mathrm{fiber}$ back along $g^{-1}$ to $x_N$; then
\eqref{eq:north-fibre-arrow} carries the result along $k$ to $y_N$; finally
$\mathcal A_A(h)$ carries that arrow forward to $Q.\mathrm{fiber}$, with
\eqref{eq:north-witness-readings} identifying its two ends. This fibre arrow
over \eqref{eq:residue-total-base-arrow} is a morphism from $P$ to $Q$ in the
ambient Grothendieck construction. Since $\iota_A$ is full it has a preimage in
the residue total $\int\mathcal R_A$, and faithfulness makes that preimage
unique. Hence $P$ and $Q$ are joined. Since they were arbitrary,
$\int\mathcal R_A$ is transitive.

Throughout the calculation $z_1$ and $z_2$ are arbitrary C3 coordinates.
Their common real value is extracted from the singleton colimit in the next
step. The comparison above rests on the relative stabilizer cancellation
through the single common source tape and north frame.
\end{proof}

\begin{theorem}[Transitivity and the singleton component of the C-residue total]\label{thm:residue-total-singleton}
\lean{ASection.sweepTransitive_on_residueSystem, ASection.residueTotal_isConnected,
ASection.residueTotal_pi0_singleton,
ASection.residueTotal_pi0_colimit_singleton,
ASection.residueInputTotalTransportRead_certified,
ASection.transportLevel_eq_sphereZero_re,
ASection.transportLevel, ASection.residueActionState_mem}
\leanok
\uses{def:A-section, def:base, def:residue-subdiagram, lem:c-residue-transitive,
lem:residue-spheres, lem:pi0-grothendieck,
thm:G2-S6, thm:winding-lift, prop:winding-signature, thm:identity}
Let $A$ be an $A$-section, and let $\int_{\mathcal B}\mathcal R_A$ be the total of its C-residue
inverse-image system, whose objects are the infinitely many C-residue zero-sphere states
transported by the A-section functor $\mathcal A_A:\mathcal B\to\mathbf{Grpd}$ from the
projective great-circle base $\mathcal B$ to the A-generated normalized value groupoids.
\begin{enumerate}[leftmargin=2.4em]
\item[(i)] $\int_{\mathcal B}\mathcal R_A$ is transitive: any two of its objects are joined by
  a morphism. Hence it is connected.
\item[(ii)] Its set of connected components is a singleton, and so is its component colimit
  $\operatorname*{colim}_{\mathcal B}(\pi_0\circ\mathcal R_A)$.
\item[(iii)] Every certified residue representative carries a real read, the real coordinate of
  its zero-sphere: $\operatorname{transportLevel}_A(n)=\re\bigl(A.\operatorname{sphereZero}(n)\bigr)$.
\end{enumerate}
The singleton in (ii) identifies the component classes of the representatives in (iii). It does
not identify their reads; see the end of the proof, Proposition~\ref{prop:read-equivalence}, and
Proposition~\ref{prop:countermodel}.
\end{theorem}
\begin{proof}
We consider a slice-preserving section $A$ of the ring $\mathcal R$ on the compactified
octonions $\OO^{*}$. Slice preservation hands us a coherently identified family of normalized
Riemann spheres, naturally organized by a functor. The
hypotheses push in the same direction. Specifying the analytic object, the theory suggests
exponential equations; and $\OO^{*}$ comes from the Cayley--Dickson construction with only one
real axis. The projective great circle that axis compactifies to is framed by $PGL(2,\RR)$: a
projective coordinate frame in which $\exp(I\theta)$ lives, and the analytic sections of
$\mathcal R$ themselves transform as M\"obius transformations --- self-transformations of the
slice-preserving section. Frames and transformations compose, invert, and act: they are
groupoids. The shape of the ring therefore determines the base $\mathcal B$ of
Definition~\ref{def:base} \textup{(}\texttt{GreatCircle.Base}\textup{)}.

Conditions C1--C4 now determine the distinguished action of the A-section. The exponential
equations lead to the commuting triangle $\pi\circ E=\exp$ of the logarithm manifold
\textup{(}Theorems~\ref{thm:slice-exp} and~\ref{thm:log-manifold}\textup{)}, giving the unique
tame lift of an exponential base \textup{(}Theorem~\ref{thm:winding-lift},
Proposition~\ref{prop:winding-signature}\textup{)}. In the matrix dictionary of
Definition~\ref{def:base}, every instant of a lift $\Gamma$ is the actual group element
\[
 D(\Gamma(t))
 =\left[C
   \begin{pmatrix}e^{\Gamma(t)}&0\\[2pt]0&1\end{pmatrix}
   C^{-1}\right].
\]
This element is also the slice function $w\mapsto D(\Gamma(t))(w)$. Hence the analytic and
symbolic readings are two faces of the same object: the exponential equation proves which matrix
acts, and matrix multiplication proves how those actions compose.

Definition~\ref{def:base} has already constructed the complete prime-indexed GPV tape beside the
matrix element it generates, including its exponential square, uniqueness, continuous real
level, winding equation, and transport to the north frame. Condition C3 supplies its Weierstrass
presentation at $N$
\textup{(}Proposition~\ref{prop:weierstrass}\textup{)}. C1 continues both presentations to the
same nonzero pole factor, whose logarithm selects $D_A$. Consequently the Euler and Weierstrass
faces are the two boundary presentations of the same matrix-built function/group element at its
fixed points $0$ and $N$. The degenerate fibre of
Lemma~\ref{lem:exp-degenerate}, with its real level and winding heights, is this action's own
bookkeeping.

Orbit--stabilizer now sweeps the element across the base. The canonical $o_b$ positions $D_A$ at
each footpoint; the uniquely forced $r_f=o_y^{-1}fo_x$ supplies the residual north action; and the
identity and multiplication laws for $r_f$ give the functor laws. The $G_2$-equivariant realization
\textup{(}Theorem~\ref{thm:G2-equiv}\textup{)} applies the same positioned M\"obius function to
the normalized value states along the continuum of Riemann spheres. The result is the production
A-section functor $\mathcal A_A=\texttt{AsectionActionDiagram\,A}$ of
Definition~\ref{def:base}; its states carry their authored inputs, positioned values, and realized
A-outputs, and its arrows are the action's actual two-legged transports.

From the functor we form its total Grothendieck construction,
$\mathcal T_A=\int_{\mathcal B}\mathcal A_A$ \textup{(}Definition~\ref{def:base}\textup{)}: an
object is a base position together with a normalized value state in its fibre, a morphism is a
base channel together with the compatible transport, and the whole is an action groupoid
formed from the distinguished element. The residue-$\CC$ zeros occur as states of this total
categorical object. The distinguished element fixes $0$ and $N$, and between them
the upper half-space; at the north frame, where the positioning factor is trivial and the
frame is the element itself, the equation $A=0$ on that half-space selects the residue states,
with Euler presenting at one end, Weierstrass at the other. The preimage of this selection
under the action --- the preimage of the total action groupoid, membership travelling by
composition in the base --- is the residue subdiagram $\mathcal R_A$ of
Definition~\ref{def:residue-subdiagram}, included in $\mathcal A_A$ by $\iota_A$. The
restriction of the total object to this preimage sits over the same base.

This restriction carries the same unique fixed $0$-to-$N$ tape for every residue zero. At the
north object the canonical positioning factor is the identity, so the frame is $D_A$ itself. The
checked north-frame and endpoint-fixation declarations identify its two fixed boundary points.
The Euler and Weierstrass expressions are the two endpoint presentations of this one action, and
the checked GPV continuation between them is unique once its initial rung is fixed. Moreover every
projective arrow has the factorization
\[
  f=o_y\,r_f\,o_x^{-1},
  \qquad r_f=\operatorname{stabilizerPart}(f),
\]
and \texttt{stabilizerPart\_unique} identifies $r_f$ as the unique north-stabilizer element in
this factorization. Thus the inverse-image witness carries the already determined $0$-to-$N$ tape
and its unique residual stabilizer part; a C3 zero index selects a state on that fixed base
channel.

\emph{The fibrewise inverse image.}
For each base object $X$, an object of $\mathcal R_A(X)$ consists of an ambient state
$x\in\mathcal A_A(X)$ together with a north residue state $x_N$ and an arrow
$g:N\to X$ satisfying
\[
  \mathcal A_A(g)(x_N)=x.                                      \tag{W}
\]
This is the object property \texttt{IsCResidueState\,A\,X}.  The fibre
$\mathcal R_A(X)$ is the full subgroupoid of $\mathcal A_A(X)$ on precisely these objects.
Its inclusion
\[
  \iota_{A,X}:\mathcal R_A(X)\hookrightarrow\mathcal A_A(X)
\]
sends a witnessed object to its ambient state.  For two such objects $x,y$, the induced map
\[
  \operatorname{Hom}_{\mathcal R_A(X)}(x,y)
  \longrightarrow
  \operatorname{Hom}_{\mathcal A_A(X)}
    (\iota_{A,X}x,\iota_{A,X}y)                       \tag{H}
\]
is a bijection: every ambient arrow between the selected states is already an arrow of the full
subgroupoid, and its ambient arrow determines it uniquely.  Surjectivity of \textup{(H)} is
fullness of $\iota_{A,X}$; injectivity is faithfulness.

\emph{Transport of the inverse image.}
Let $f:X\to Y$.  If $x$ has witness \textup{(W)}, then $\mathcal A_A(f)(x)$ has the same north state
$x_N$ and the composite witness $g\fatsemi f:N\to Y$, because
\[
\begin{aligned}
  \mathcal A_A(g\fatsemi f)(x_N)
    &=\mathcal A_A(f)\bigl(\mathcal A_A(g)(x_N)\bigr)\\
    &=\mathcal A_A(f)(x).
\end{aligned}                                         \tag{C}
\]
This is the calculation recorded by \texttt{cResidue\_lands}.  It defines the restricted
transport
\[
  \mathcal R_A(f):\mathcal R_A(X)\longrightarrow\mathcal R_A(Y):
\]
on objects it uses the composite witness in \textup{(C)}, and on arrows it applies the ambient
transport $\mathcal A_A(f)$.  Composition in the projective base therefore supplies closure of the
inverse-image system under every A-section transport.

\emph{The natural inclusion.}
The fibrewise inclusions assemble over $\mathcal B$ in the square
\[
\begin{tikzcd}[column sep=huge,row sep=large]
\mathcal R_A(X) \arrow[r,"\mathcal R_A(f)"] \arrow[d,hook,"\iota_{A,X}"'] &
\mathcal R_A(Y) \arrow[d,hook,"\iota_{A,Y}"]\\
\mathcal A_A(X) \arrow[r,"\mathcal A_A(f)"'] & \mathcal A_A(Y).
\end{tikzcd}                                           \tag{Nat}
\]
For an object $x$, both routes in \textup{(Nat)} have underlying state $\mathcal A_A(f)(x)$; for an arrow
$\alpha$, both routes have underlying arrow $\mathcal A_A(f)(\alpha)$.  The upper route additionally
carries the residue witness $g\fatsemi f$ established in \textup{(C)}.  Hence
\[
  \mathcal R_A(f)\fatsemi \iota_{A,Y}
  =\iota_{A,X}\fatsemi \mathcal A_A(f).                    \tag{Inc}
\]
and these equations for every $f$ are the naturality of
$\iota_A:\mathcal R_A\Rightarrow\mathcal A_A$.

The preimage now assembles over the base into a total of its own. Its objects are the positioned
squares of the C-residue image and its morphisms are the value transports of $\mathcal T_A$
between them. Its inclusion is the total reading of $\iota_A$,
\[
\begin{tikzcd}[column sep=huge]
\int_{\mathcal B}\mathcal R_A \arrow[r, hook, "\int\iota_A"] & \int_{\mathcal B}\mathcal A_A=\mathcal T_A
\end{tikzcd}
\]
\textup{(}\texttt{Asection\-CResidue\-Inclusion\-Total} in the Lean source, the Grothendieck
functoriality of $\iota_A$ applied to it\textup{)}. On objects it sends $(X,x)$ to
$(X,\iota_{A,X}x)$ and leaves the base position unchanged; on a Grothendieck morphism it leaves
the base arrow unchanged and applies the fibre inclusion only to the compatible transport.
Consequently equality after inclusion forces equality of the base legs and, by fibrewise
faithfulness, equality of the transport legs: the total inclusion is faithful. Conversely, an
ambient total morphism between two objects in the image already has a base leg $f$ and a fibre
morphism between selected target states; fibrewise fullness gives its preimage in
$\mathcal R_A$, with the same $f$. Hence the total inclusion is full
\textup{(}\texttt{Asection\-CResidue\-Inclusion\-Total\-\_full} and
\texttt{Asection\-CResidue\-Inclusion\-Total\-\_faithful}\textup{)}. It is therefore an
isomorphism onto its image. The C-residue total is that image itself: it consists exactly of the
selected objects together with the fixed tape and all ambient total morphisms between them.

Lemma~\ref{lem:c-residue-transitive} proves, in this exact inverse-image groupoid, that any
two objects are joined by one Grothendieck morphism: the relative
orbit--stabilizer face is its base component and the $G_2$ transport is its fibre component.
Thus $\int_{\mathcal B}\mathcal R_A$ is connected.  This is a groupoid-register statement:
at this stage the established conclusion is connectedness; equality of real coordinates and the
common centre arise from the colimit readout below.

By Remark~8.3.5 of \cite{Riehl}, a category is nonempty and connected if and only if its
$\pi_0$ is the singleton set. Condition C4 supplies nonemptiness --- infinitely many
residue-$\CC$ zero-sphere states populate the preimage --- and connectedness was just
established. Therefore, by Lemma~\ref{lem:pi0-grothendieck}
\textup{(}\texttt{pi0GrothendieckEquiv}\textup{)},
\[
\pi_0\Bigl(\int_{\mathcal B}\mathcal R_A\Bigr)\;\cong\;\operatorname*{colim}_{\mathcal B}\,(\pi_0\circ\mathcal R_A)
\]
collapses to a singleton $\kappa$. This proves (i) and (ii).

For (iii): the $n$-th certified residue representative is the state of the $n$-th enumerated
zero-sphere at the north frame, carried along the fixed tape; its positioned coordinate is the
sphere parameter $q_n=A.\operatorname{sphereZero}(n)$, and the read
$\operatorname{transportLevel}_A(n)$ is by definition the real part of that coordinate
\textup{(}\texttt{residueInputTotalTransportRead\_certified},
\texttt{transportLevel\_eq\_sphereZero\_re}\textup{)}.

This is where the argument stops. The singleton $\kappa$ identifies the component classes of any
two certified representatives. It does not identify their reads: two objects of one component
are joined by a zigzag of transports of $\int_{\mathcal B}\mathcal R_A$, and no cocone, descent
law, or invariance law for the read along those transports is established. Equality of the reads
is the concentricity statement itself \textup{(}Proposition~\ref{prop:read-equivalence}\textup{)},
and it fails for the A-section of Proposition~\ref{prop:countermodel}, so no such law can hold
at this generality. Earlier versions of this theorem concluded from $\kappa$ that all reads
equal one value $c$ and that the zero-spheres are concentric; that step was never accepted by
Lean, and its conclusion is false. It was withdrawn on 2026-09-02.
\end{proof}

\begin{proposition}[The read equivalence]\label{prop:read-equivalence}
\lean{ConcentricASection, PairwiseTransportLevel,
concentricASection_iff_pairwiseTransportLevel}
\leanok
\uses{def:A-section, def:residue-subdiagram, thm:residue-total-singleton}
For an A-section $A$ the following are equivalent:
\begin{enumerate}[leftmargin=2.4em]
\item[(a)] the enumerated residue-$\CC$ zero-spheres $\{q_n\}$ share one real centre: there is
  $c\in\RR$ with $\re q_n=c$ for every $n$;
\item[(b)] any two certified residue representatives carry the same real read:
  $\operatorname{transportLevel}_A(n)=\operatorname{transportLevel}_A(m)$ for all $n,m$.
\end{enumerate}
\end{proposition}
\begin{proof}
By Theorem~\ref{thm:residue-total-singleton}(iii), $\operatorname{transportLevel}_A(n)=\re q_n$
for every $n$. If (a) holds, both reads in (b) equal $c$. If (b) holds, every $\re q_n$ equals
$\re q_0$, so $c=\re q_0$ witnesses (a).
\end{proof}

\begin{proposition}[An A-section whose zero-spheres are not concentric]\label{prop:countermodel}
\lean{SpecCandidate.candidateSection, SpecCandidate.candidateSection_not_concentric,
SpecCandidate.current_ASection_concentricity_type_false,
SpecCandidate.candidateSection_not_pairwiseTransportLevel}
\leanok
\uses{def:A-section, cor:zeta-section, thm:zero-spheres, prop:weierstrass, prop:read-equivalence}
Let $A=(q^{2}+1)\,\zetaO$, the octonionic zeta multiplied by the slice-regular primary factor of
the zero-sphere through $i$. Then $A$ satisfies \textup{(C1)--(C4)}: the pole is at $1$; the
Euler family of $\zetaO$, indexed by the primes, is augmented by a second infinite family,
indexed by $\mathbb{N}$, of geometric scalings $\tfrac12\,2^{-n}\,\ell$ of one slice-preserving
logarithm $\ell$ of $q^{2}+1$, each zero-free and slice regular on the half-space $\re q>1$ and
summing there to $\ell$; the real Weierstrass factor and the entire factor are those of
$\zetaO$; and the residue-$\CC$ zero-spheres are enumerated by $q_0=i$, $q_{n+1}=\rho_n$, where
$\{\rho_n\}$ enumerates the upper-half-plane zeros of $\zeta$. Its zero-spheres are not
concentric: $\re q_0=0$, while $0<\re\rho_n<1$ for every $n$. Consequently the statement
``the residue-$\CC$ zero-spheres of every A-section are concentric'' is false, and by
Proposition~\ref{prop:read-equivalence} the certified reads of $A$ are not pairwise equal.
\end{proposition}
\begin{proof}
Every clause of Definition~\ref{def:A-section}, in the form formalized as the structure
\texttt{ASection}, is discharged in \texttt{Concentricity/SpecCandidate.lean}: the factor
$q^{2}+1$ is slice preserving and entire with the single zero-sphere through $i$; the
conjugate-symmetrized principal logarithm $\ell$ of $q^{2}+1$ agrees with the principal logarithm
on $\re q>1$, where it is zero-free and analytic, so its geometric scalings are zero-free, analytic,
and summable alongside the Euler logarithms of $\zeta$, with a local-normal majorant; the product
$(q-1)A$ has the stated full factorization,
with the convergence of Proposition~\ref{prop:weierstrass}; and the enumeration is infinite
because the zeta enumeration is. The real parts are $0$ for $q_0$ and, for the inherited zeros,
strictly between $0$ and $1$ by the open-strip theorem (Theorem~\ref{thm:riemann}, as formalized
in \texttt{nontrivialZero\_re\_mem\_Ioo}). The point $i$ is not a zero of $\zeta$, so it is a
genuine additional zero-sphere of $A$.
\end{proof}

\begin{remark}[Translation to the classical framework]\label{rmk:nontrivial}
Under the residue-$\CC$/classically-nontrivial-zero identification
\textup{(}Theorems~\ref{thm:zero-equivalence} and~\ref{thm:zero-spheres}\textup{)}, the
populated residue-$\CC$ zero spheres of an A-section are its classically nontrivial zero spheres,
infinite and pairwise disjoint by \textup{C4} and Lemma~\ref{lem:residue-spheres}. Whether they
share a common centre is condition (a) of Proposition~\ref{prop:read-equivalence}; it is not
supplied by Theorem~\ref{thm:residue-total-singleton}, and it fails for the A-section of
Proposition~\ref{prop:countermodel}. Earlier versions of this paper stated an unconditional
corollary here, carrying a common centre supplied by the former universal theorem; it was
withdrawn on 2026-09-02.
\end{remark}

% The superseded great-circle/F-based proof frame was removed from the active
% master on 2026-07-17. It remains recoverable from git history and is not an
% operative proof or expository alternative.

\section{Corollaries}\label{sec:verification}

Two results close the argument. The first checks that the octonionic zeta is an A-section,
verifying its four conditions C1--C4 one by one against the classical facts of Part~1, so that
Theorem~\ref{thm:residue-total-singleton} applies to it. The second says exactly what
concentricity means for it: its residue-$\CC$ zero-spheres share one real centre if and only if
the Riemann Hypothesis holds, the same is true of the certified reads, and under the Riemann
Hypothesis the common centre is $\tfrac12$, pinned by the functional equation of $\xi$ through
Theorem~\ref{thm:rh-equiv}.

\begin{corollary}[$\zetaO$ is an $A$-section]\label{cor:zeta-section}
\lean{zetaSection}
\uses{def:A-section, thm:zeta-in-R, thm:riemann, thm:euler, cor:hadamard-infinitude, prop:weierstrass, thm:zero-spheres}
$\zetaO$ is an $A$-section \textup{(Definition~\ref{def:A-section})}: a section of
$\mathcal R$ satisfying \textup{(C1)--(C4)}.
\end{corollary}
\begin{proof}
$\zetaO\in\mathcal R$ by slice-preservation (Theorem~\ref{thm:zeta-in-R}). The verification of
\textup{(C1)--(C4)} cites the classical facts of the Riemann zeta function (Part~1) and the
slice-regular zero-geometry that translates them:
\begin{itemize}[leftmargin=2em]
\item \textbf{(C1)} \emph{Meromorphic continuation; simple pole.} Classical fact: $\zeta$ is
  meromorphically continued with a single simple pole at $s=1$ (Theorem~\ref{thm:riemann});
  the value of $\zetaO$ there is the point at infinity, $\zetaO(1)=\infty=N$
  (Definition~\ref{def:zeta_O}).
\item \textbf{(C2)} \emph{Infinite Euler product.} Classical fact: $\zeta=\prod_p(1-p^{-s})^{-1}
  =\exp(\sum_p\ell_p)$ on $\Re s>1$, the family indexed by the infinitely many primes
  (Theorem~\ref{thm:euler}).
\item \textbf{(C3)} \emph{Infinite Weierstrass factorization.} Supplied by the
  \emph{slice-regular} Weierstrass factorization (Proposition~\ref{prop:weierstrass};
  \cite[Prop.~3.1, Thm.~3.2]{AdFslice}, \cite{GV}):
  $(q-1)\,\zetaO=q^{m}Re^{g}\prod_n\mathcal E(\cdot;q_n)$ \textup{(}pole factor at $p_0=1$;
  classically, Hadamard factors $(s-1)\zeta(s)$\textup{)}, where $m=0$
  \textup{(}$\zeta(0)=-\tfrac12\neq0$\textup{)}, $R$ is the product over the nonzero
  residue-$\RR$ zeros, the classically trivial zeros $-2,-4,\dots$, honest real zeros of
  $\zetaO$, and $\{q_n\}$ enumerates its residue-$\CC$ zero-spheres.
\item \textbf{(C4)} \emph{Infinitely many residue-$\CC$ zeros.} Classical fact: $\zeta$ has
  infinitely many nontrivial zeros (Corollary~\ref{cor:hadamard-infinitude}; also Hardy,
  Theorem~\ref{thm:hardy}). By the residue dictionary below these are exactly the residue-$\CC$
  zeros, so $\{q_n\}$ is infinite.
\end{itemize}
After the \textup{(C1)--(C4)} verification, the downstream centre-pinning uses the
\emph{functional equation} $\xi(s)=\xi(1-s)$,
$\xi(\bar s)=\overline{\xi(s)}$ (Theorem~\ref{thm:riemann}), used in Theorem~\ref{thm:zeta-concentricity-rh}. The
\emph{residue dictionary}, a closed zero of $\zetaO$ is a real point (residue field $\RR$, the
classically \emph{trivial} zeros) or a $6$-sphere $S_{(\sigma,\gamma)}$ (residue field $\CC$, the
classically \emph{nontrivial} zeros), is the slice-regular zero-geometry of
Theorem~\ref{thm:zero-spheres} and \cite{AdF}; it is what makes ``residue-$\CC$'' and
``nontrivial'' the same notion. Thus \textup{(C1)--(C4)} hold.
\end{proof}

\begin{theorem}[Concentricity of $\zetaO$ is the Riemann Hypothesis]\label{thm:zeta-concentricity-rh}
\lean{zetaSection_concentric_iff_RH, zetaSection_pairwiseTransportLevel_iff_RH,
zetaSection_sphereZero_re_eq_half_of_RH, zetaSection_transportLevel_eq_half_of_RH,
zeta_criticalLine_zeros_infinite_of_RH}
\leanok
\uses{cor:zeta-section, thm:rh-equiv, prop:read-equivalence, thm:residue-total-singleton,
thm:zero-spheres, cor:hadamard-infinitude}
For the A-section $\zetaO$ of Corollary~\ref{cor:zeta-section}:
\begin{enumerate}[leftmargin=2.4em]
\item[(a)] the enumerated residue-$\CC$ zero-spheres $\{q_n\}$ share one real centre if and only
  if the Riemann Hypothesis holds;
\item[(b)] the certified reads $\operatorname{transportLevel}_{\zetaO}(n)$ are pairwise equal if
  and only if the Riemann Hypothesis holds;
\item[(c)] under the Riemann Hypothesis, every $q_n$ has real part $\tfrac12$, every read
  $\operatorname{transportLevel}_{\zetaO}(n)$ equals $\tfrac12$, and infinitely many zeros of
  $\zeta$ lie on the line $\Re s=\tfrac12$.
\end{enumerate}
\end{theorem}
\begin{proof}
(a) The enumeration $\{q_n\}$ of Corollary~\ref{cor:zeta-section} is complete: every zero of
$\zeta$ in the upper half-plane is some $q_n$
\textup{(}\texttt{zetaSphereZero\_surjective}\textup{)}, and every $q_n$ is such a zero. Hence a
common real centre for $\{q_n\}$ is the same as a common real part for all upper-half-plane
zeros of $\zeta$, which is equivalent to the Riemann Hypothesis by Theorem~\ref{thm:rh-equiv}.
(b) follows from (a) and Proposition~\ref{prop:read-equivalence}. (c) Under the Riemann
Hypothesis the common centre of Theorem~\ref{thm:rh-equiv} is $\tfrac12$; it is $\re q_n$ and,
by Theorem~\ref{thm:residue-total-singleton}(iii), the read. The infinitude clause is
Corollary~\ref{cor:hadamard-infinitude} restricted to the line.
\end{proof}

\begin{remark}[What this theorem is]\label{rmk:reframing}
Theorem~\ref{thm:zeta-concentricity-rh} is a geometric and categorical reformulation of the
Riemann Hypothesis, formalized without placeholders. It is not a proof of the Riemann Hypothesis,
and nothing in this paper or its formalization proves the Riemann Hypothesis. Clause (b) shows
that the pairwise read equality which the singleton component of
Theorem~\ref{thm:residue-total-singleton} does not supply is, for $\zetaO$, of exactly the
strength of the Riemann Hypothesis: a categorical proof of it would be a proof of the Riemann
Hypothesis. Earlier versions of this paper stated the Riemann Hypothesis here as an
unconditional corollary of a universal concentricity theorem; that theorem is false
\textup{(}Proposition~\ref{prop:countermodel}\textup{)} and the corollary was withdrawn on
2026-09-02. Hardy's theorem (Theorem~\ref{thm:hardy}) remains a cited classical fact and is not
formalized.
\end{remark}

A word on why the argument takes the shape it does. No one can hold an $8$-dimensional graph in the
mind's eye, or draw the continuum of Riemann spheres the section lives on. The epigraph names the
way through: no single viewpoint sees the whole, but when the viewpoints are multiplied and
conjugated, a vision appears that none of them held alone. Each slice direction is one such
viewpoint; the value-bearing A-section functor gathers all of them over the base, the total object holds them
together, and the colimit conjugates them into a single component. Whether that component's
real reads agree is the concentricity question itself: the component we can compute, the
agreement we cannot, and for the octonionic zeta the agreement is the Riemann Hypothesis.

\begin{remark}[The image span of a fully faithful inclusion]\label{rmk:image-span}
For any fully faithful functor $F:\mathcal C\to\mathcal D$ there is a canonical roof through
the image,
\[
\begin{tikzcd}[column sep=large]
& \operatorname{Im}(F) \arrow[dl, "F^{-1}"'] \arrow[dr, hook, "j"] & \\
\mathcal C & & \mathcal D
\end{tikzcd}
\]
the backward leg being an isomorphism of categories onto the image and the forward leg the
canonical inclusion. It is a general fact, available here for $\iota_A$ and for
$\int\iota_A$, and it is the shape a length-two zigzag argument would use.
Theorem~\ref{thm:residue-total-singleton} obtains connectedness directly from transitivity at the total,
by a single element. The image span records the corresponding length-two alternative.
\end{remark}

\begin{remark}[Two alternative routes]\label{rmk:thomason}
Two other proofs of Theorem~\ref{thm:residue-total-singleton} are available. The first runs over the
base. The restriction sits over $\mathcal B$ through its two projections,
\[
\begin{tikzcd}[column sep=large]
\int_{\mathcal B}\mathcal R_A \arrow[rr, hook, "\int\iota_A"] \arrow[dr, "\pi_{\mathcal R}"'] & &
\int_{\mathcal B}\mathcal A_A \arrow[dl, "\pi_{\mathcal A}"] \\
& \mathcal B &
\end{tikzcd}
\]
and one may connect residue states footpoint by footpoint: zigzags in the connected base,
lifted through the fibres. That route travels through $\mathcal B$ and reassembles the value data
leg by leg. The second is Thomason's homotopy colimit theorem \cite{Thomason79}: the
classifying space of a Grothendieck construction is homotopy equivalent to the homotopy
colimit of the classifying spaces of its fibres, and reading connected components through that
equivalence yields the same collapse of $\pi_0\bigl(\int_{\mathcal B}\mathcal R_A\bigr)$. The
proof given here stays at the level of categories, where Lemma~\ref{lem:pi0-grothendieck}
suffices: the base route assembles what the action already supplies whole, the Thomason route
is the topological shadow of the same computation, and both remain expository routes outside
the Lean development.
\end{remark}

%=====================================================================
\begin{thebibliography}{99}
%=====================================================================
\bibitem{AdF} A.~Altavilla and C.~De Fabritiis, \emph{$\ast$-logarithm for slice regular
functions}, Rend.\ Lincei Mat.\ Appl.\ (arXiv:2106.04227); the slice-regular exponential
$\exp_{*}$ (Def.~2.14, and Rem.~2.23 for the slice-preserving formula
$\exp_{*}(f)=e^{f_0}(\cos f_1+I\sin f_1)$) and the $\ast$-logarithm on never-vanishing
functions (Def.~2.24, Thm.~6.17).
\bibitem{AdFslice} A.~Altavilla and C.~De Fabritiis, \emph{$s$-regular functions which
preserve a complex slice}, Ann.\ Mat.\ Pura Appl.\ (arXiv:1801.01318); Prop.~3.1 and
Thm.~3.2 factor a slice-regular function into a factor carrying exactly its (real and
spherical) zeros and a never-vanishing factor, including the case of infinitely many zeros,
the citable form of C3.
\bibitem{Baez02} J.~C.~Baez, \emph{The octonions}, Bull.\ Amer.\ Math.\ Soc.\
\textbf{39} (2002), 145--205.
\bibitem{BisiWinkelmann} C.~Bisi and J.~Winkelmann, \emph{Invariants and Automorphisms for slice
regular functions: the octonionic case}, arXiv:2411.16762 (2024); companion arXiv:2411.15896 in
J.~Noncommutative Geometry (EMS). Supplies the $I$-independent lift $\phi_I\circ F=A\circ\phi_I$
(the slice projection) and, in \S3.2/\S3.7/\S14, the axially-symmetric-domain, slice-preserving, and
central-divisor notions.
\bibitem{CSS12} F.~Colombo, I.~Sabadini, and D.~C.~Struppa, \emph{Noncommutative
Functional Calculus: Theory and Applications of Slice Hyperholomorphic Functions},
Progress in Mathematics \textbf{289}, Birkh\"auser, 2011.
\bibitem{GP11} R.~Ghiloni and A.~Perotti, \emph{Slice regular functions on real
alternative algebras}, Adv.\ Math.\ \textbf{226} (2011), 1662--1691.
\bibitem{SeriesExp} R.~Ghiloni, A.~Perotti, and C.~Stoppato, \emph{Singularities of slice regular
functions over real alternative $\ast$-algebras}, Adv.\ Math.\ \textbf{305} (2017), 1085--1130
(slice-function and stem-function formalism, \S1--\S2; the slice-preserving
$\Leftrightarrow$ intrinsic $\Leftrightarrow$ $\RR$-valued-stem characterization used in
Definition~\ref{def:slice-squares}; \S11: slice semiregular functions, Def.~11.1, Thm.~11.3).
\bibitem{GS07} G.~Gentili and D.~C.~Struppa, \emph{A new theory of regular functions
of a quaternionic variable}, Adv.\ Math.\ \textbf{216} (2007), 279--301.
\bibitem{GSS13} G.~Gentili, C.~Stoppato, and D.~C.~Struppa, \emph{Regular Functions of
a Quaternionic Variable}, Springer Monographs in Mathematics, Springer, 2013.
\bibitem{GV} G.~Gentili and I.~Vignozzi, \emph{The Weierstrass factorization theorem
for slice regular functions over the quaternions}, Ann.\ Global Anal.\ Geom.\
\textbf{40} (2011), 435--466.
\bibitem{Hardy14} G.~H.~Hardy, \emph{Sur les z\'eros de la fonction $\zeta(s)$ de
Riemann}, C.\ R.\ Acad.\ Sci.\ Paris \textbf{158} (1914), 1012--1014.
\bibitem{Munkres00} J.~R.~Munkres, \emph{Topology}, 2nd ed., Prentice Hall, 2000.
\bibitem{Riemann1859} B.~Riemann, \emph{\"Uber die Anzahl der Primzahlen unter einer
gegebenen Gr\"osse}, Monatsber.\ Berlin.\ Akad.\ (1859).
\bibitem{Schafer66} R.~D.~Schafer, \emph{An Introduction to Nonassociative Algebras},
Academic Press, 1966.
\bibitem{Titchmarsh86} E.~C.~Titchmarsh (rev.\ D.~R.~Heath-Brown), \emph{The Theory of
the Riemann Zeta-Function}, 2nd ed., Oxford University Press, 1986.
\bibitem{Wang} X.~Wang, \emph{On geometric aspects of quaternionic and octonionic
slice regular functions}, arXiv:1512.01414v3.
\bibitem{VS} G.~Gentili, J.~Prezelj, and F.~Vlacci, \emph{Slice conformality and
Riemann manifolds on quaternions and octonions}, Math.\ Z.\ \textbf{302} (2022),
971--994 (arXiv:2107.07892v2).
\bibitem{Stacks} The Stacks Project Authors, \emph{The Stacks Project},
\url{https://stacks.math.columbia.edu}, 2024.
\bibitem{Mathlib} The mathlib Community, \emph{The Lean mathematical library}
(\texttt{CategoryTheory.ActionCategory}, \texttt{CategoryTheory.ConnectedComponents}, \texttt{CategoryTheory.Comma}),
\url{https://leanprover-community.github.io/mathlib4_docs/}.
\bibitem{GJ} P.~G.~Goerss and J.~F.~Jardine, \emph{Simplicial Homotopy Theory}, Modern
Birkh\"auser Classics, Birkh\"auser, 2009 (Ch.~I: simplicial sets and the nerve; Ch.~IV: the
bisimplicial engine for homotopy colimits).
\bibitem{GPVwind} G.~Gentili, J.~Prezelj, and F.~Vlacci, \emph{On a continuation of
quaternionic and octonionic logarithm along curves and the winding number},
J.~Math.\ Anal.\ Appl.\ \textbf{536} (2024), no.~1, Paper No.~128219,
DOI 10.1016/j.jmaa.2024.128219 (arXiv:2307.14047) (Rem.~2.1: the direction $I(q)$ has no
continuous extension to $\RR$; Def.~4.7: tame path = unique companion (Def.~4.20 for maps;
Def.~5.2: tame/semi-tame at an obstruction parameter); Def.~5.11: the loop lift
$\mathrm{pr}_1\circ\Gamma=\gamma\circ\exp$; Cor.~5.13: a lift of a loop exists iff the
signature lies in $\{0,-1\}$ on each obstruction interval, and is then a loop; Cor.~5.21:
winding $=|\sigma^c|/2$).
\bibitem{Quillen73} D.~Quillen, \emph{Higher algebraic $K$-theory: I}, in \emph{Algebraic
$K$-theory, I}, Lecture Notes in Math.\ \textbf{341}, Springer, 1973, 85--147 (Theorem~A:
a functor whose comma categories are contractible is final and induces an equivalence of
classifying spaces).
\bibitem{Riehl} E.~Riehl, \emph{Categorical Homotopy Theory}, New Mathematical Monographs \textbf{24},
Cambridge University Press, 2014 (the nerve, classifying spaces, the two-sided bar construction, and
homotopy colimits; Ch.~1, 4--6).
\bibitem{RiehlContext} E.~Riehl, \emph{Category Theory in Context}, Aurora: Dover Modern Math
Originals, Dover Publications, 2016 (action groupoids and their orbits; Ex.~1.5.19).
\bibitem{Thomason79} R.~W.~Thomason, \emph{Homotopy colimits in the category of small
categories}, Math.\ Proc.\ Cambridge Philos.\ Soc.\ \textbf{85} (1979), 91--109 (Thm.~1.2:
the classifying space of the Grothendieck construction as a homotopy colimit).
\bibitem{Vistoli} A.~Vistoli, \emph{Notes on Grothendieck topologies, fibered categories and descent
theory}, in \emph{Fundamental Algebraic Geometry}, Math.\ Surveys Monogr.\ \textbf{123}, AMS, 2005
(arXiv:math/0412512); \S3.1 for the Grothendieck construction and categories fibered in groupoids.
\end{thebibliography}

\end{document}
